\section{AMC 10 10A 2021}
        \begin{enumerate}
        	\item (\textit{AMC 10 10A 2021}) What is the value of
 \[(2^2-2)-(3^2-3)+(4^2-4)?\]$\textbf{(A)} ~1 \qquad\textbf{(B)} ~2 \qquad\textbf{(C)} ~5 \qquad\textbf{(D)} ~8 \qquad\textbf{(E)} ~12$

 \textbf{Answer: }D



\textbf{Solution 1}

$(4-2)-(9-3)+(16-4)=2-6+12=8.$ This corresponds to $\boxed{\textbf{(D) } 8}.$

 -happykeeper

 



\textbf{Solution 2}

We have
 \begin{align*} \left(2^2-2\right)-\left(3^2-3\right)+\left(4^2-4\right) &= 2(2-1)-3(3-1)+4(4-1) \\ &= 2(1)-3(2)+4(3) \\ &= 2-6+12 \\ &= \boxed{\textbf{(D) } 8}. \end{align*}
~MRENTHUSIASM

 



\textbf{Solution 3 (Overkill: Just for Fun)}

We have
 \begin{align*} (2^2-2)-(3^2-3)+(4^2-4) &= 2^2+4^2-3^2-2+3-4 \\ &=2^2+(4-3)(4+3)-3 \\ &=2^2+7-3=2^2+4 \\ &=4\cdot 2 \\ &=\boxed{\textbf{(D) } 8}. \end{align*}
-\href{https://artofproblemsolving.com/wiki/index.php?title=User:PureSwag}{PureSwag}

 



\textbf{Solution 4 (When you have too much time)}

Using the difference of squares, we have 
 \[(2-\sqrt{2})(2+\sqrt{2}) - (3-\sqrt{3})(3+\sqrt{3}) + (4-2)(4+2).\]Knowing that $\sqrt{2} \approx 1.41$ and $\sqrt{3} \approx 1.73,$ we get 
 \[(2-\sqrt{2})(2+\sqrt{2}) - (3-\sqrt{3})(3+\sqrt{3}) + (4-2)(4+2) \approx 0.59\cdot 3.41 -1.26\cdot 4.73 + 2 \cdot 6 =8.0521,\] which is closest to $\boxed{\textbf{(D) } 8}.$



	\item (\textit{AMC 10 10A 2021}) Portia's high school has $3$ times as many students as Lara's high school. The two high schools have a total of $2600$ students. How many students does Portia's high school have?


 $\textbf{(A)} ~600 \qquad\textbf{(B)} ~650 \qquad\textbf{(C)} ~1950 \qquad\textbf{(D)} ~2000\qquad\textbf{(E)} ~2050$

 \textbf{Answer: }C



\textbf{Solution 1 (Two Variables)}

The following system of equations can be formed with $P$ representing the number of students in Portia's high school and $L$ representing the number of students in Lara's high school:
 \begin{align*} P&=3L, \\ P+L&=2600. \end{align*}
Substituting $P=3L$ gives $4L=2600.$ Solving for $L$ gives $L=650.$ Since we need to find $P,$ we multiply $650$ by $3$ to get $P=\boxed{\textbf{(C)} ~1950}.$




\textbf{Solution 2 (One Variable)}

Suppose Lara's high school has $x$ students, so Portia's high school has $3x$ students. We have $x+3x=2600,$ or $4x=2600.$ The answer is 
 \[3x=2600\cdot\frac 34=650\cdot3=\boxed{\textbf{(C)} ~1950}.\]




\textbf{Solution 3 (Arithmetic)}

Clearly, $2600$ is $4$ times the number of students in Lara's high school. Therefore, Lara's high school has $2600\div4=650$ students, and Portia's high school has $650\cdot3=\boxed{\textbf{(C)} ~1950}$ students.

 



\textbf{Solution 4 (Observations)}

The number of students in Portia's high school must be a multiple of $3.$ This eliminates $\textbf{(B)},\textbf{(D)},$ and $\textbf{(E)}.$ Since $\textbf{(A)}$ is too small (as it is clear that $600+\frac{600}{3}<2600$), we are left with $\boxed{\textbf{(C)} ~1950}.$



	\item (\textit{AMC 10 10A 2021}) The sum of two natural numbers is $17{,}402$. One of the two numbers is divisible by $10$. If the units digit of that number is erased, the other number is obtained. What is the difference of these two numbers?


 $\textbf{(A)} ~10{,}272\qquad\textbf{(B)} ~11{,}700\qquad\textbf{(C)} ~13{,}362\qquad\textbf{(D)} ~14{,}238\qquad\textbf{(E)} ~15{,}426$

 \textbf{Answer: }D



\textbf{Solution 1 (Algebra)}

The units digit of a multiple of $10$ will always be $0$. We add a $0$ whenever we multiply by $10$. So, removing the units digit is equal to dividing by $10$.

 Let the smaller number (the one we get after removing the units digit) be $a$. This means the bigger number would be $10a$.

 We know the sum is $10a+a = 11a$ so $11a=17402$. So $a=1582$. The difference is $10a-a = 9a$. So, the answer is $9(1582) = \boxed{\textbf{(D)} ~14{,}238}$.

 



\textbf{Solution 2 (Arithmetic)}

Since the unit's place of a multiple of $10$ is $0$, the other integer must end with a $2$, for both integers sum up to a number ending in a $2$. Thus, the unit's place of the difference must be $10-2=8$, and the only answer choice that ends with an $8$ is $\boxed{\textbf{(D)} ~14{,}238}$.

 Another quick solution is to realize that the sum represents a number $n$ added to $10n$. The difference is $9n$, which is $\frac{9}{11}$ of the given sum.

 



\textbf{Solution 3 (Vertical Addition and Logic)}

Let the larger number be $\underline{AB{,}CD0}.$ It follows that the smaller number is $\underline{A{,}BCD}.$ Adding vertically, we have 
 \[\begin{array}{cccccc}   & A & B & C & D & 0 \\ +\quad &   & A & B & C & D \\ \hline   &   &   &   &   &   \\ [-2.5ex]   & 1 & 7 & 4 & 0 & 2 \\ \end{array}\]Working from right to left, we get 
 \[D=2\implies C=8 \implies B=5 \implies A=1.\]The larger number is $15{,}820$ and the smaller number is $1{,}582.$ Their difference is $15{,}820-1{,}582=\boxed{\textbf{(D)} ~14{,}238}.$




\textbf{Solution 4 (Logic)}

We know that the larger number has a units digit of $0$ since it is divisible by 10. If $D$ is the ten's digit of the larger number, then $D$ is the units digit of the smaller number. Since the sum of the natural numbers has a unit's digit of $2$, $D=2$. 

 The units digit of the larger number is $0$ and the units digit of the smaller number is $2$, so the positive difference between the numbers is 8. There is only one answer choice that has this units digit, and that is $\boxed{\textbf{(D)} ~14{,}238}.$

 (Similar to MRENTHUSIASM's solution)

 


	\item (\textit{AMC 10 10A 2021}) A cart rolls down a hill, travelling $5$ inches the first second and accelerating so that during each successive $1$-second time interval, it travels $7$ inches more than during the previous $1$-second interval. The cart takes $30$ seconds to reach the bottom of the hill. How far, in inches, does it travel?


 $\textbf{(A)} ~215 \qquad\textbf{(B)} ~360\qquad\textbf{(C)} ~2992\qquad\textbf{(D)} ~3195\qquad\textbf{(E)} ~3242$

 \textbf{Answer: }D



\textbf{Solution 1 (Arithmetic Series)}

Since 
 \[\mathrm{Distance}=\mathrm{Speed}\cdot\mathrm{Time},\] we seek the sum 
 \[5\cdot1+12\cdot1+19\cdot1+26\cdot1+\cdots=5+12+19+26+\cdots,\] in which there are $30$ terms.

 The last term is $5+7\cdot(30-1)=208.$ Therefore, the requested sum is 
 \[5+12+19+26+\cdots+208=\frac{5+208}{2}\cdot30=\boxed{\textbf{(D)} ~3195}.\] Recall that to find the sum of an arithmetic series, we take the average of the first and last terms and then multiply by the number of terms: 
 \[\mathrm{Sum}=\frac{\mathrm{First}+\mathrm{Last}}{2}\cdot\mathrm{Count}.\] ~MRENTHUSIASM

 



\textbf{Solution 2 (Arithmetic Series)}

The distance (in inches) traveled within each $1$-second interval is: 
 \[5,5+1(7),5+2(7), \dots , 5+29(7).\]This is an arithmetic sequence so the total distance travelled, found by summing them up is:
 \[\text{number of terms} \cdot \text{average of terms} = \text{number of terms} \cdot \dfrac{\text{first term}+\text{last term}}{2}.\]Or, 
 \[30 \cdot \dfrac{5+5+29(7)}{2} = 15 \cdot 213 = \boxed{\textbf{(D)} ~3195}.\]~BakedPotato66

 



\textbf{Solution 3 (Answer Choices and Modular Arithmetic)}

From the $30$-term sum 
 \[5+12+19+26+\cdots\] in Solution 1, taking modulo $10$ gives 
 \[5+12+19+26+\cdots \equiv 3\cdot(5+2+9+6+3+0+7+4+1+8) = 3\cdot45\equiv5 \pmod{10}.\] The only answer choices congruent to $5$ modulo $10$ are $\textbf{(A)}$ and $\textbf{(D)}.$ By a quick estimation, $\textbf{(A)}$ is too small, leaving us with $\boxed{\textbf{(D)} ~3195}.$




\textbf{Solution 4 (Motion With Constant Acceleration)}

This problem can be solved by physics method. This method is perhaps the quickest too and shows the beauty of the problem. The average speed increases $7 \ \text{in/s}$ per second. So, the acceleration $a=7 \ \text{in/s\textsuperscript{2}}.$ The average speed of the first second is $5 \ \text{in/s}.$ We can know the initial velocity $v_0=5-0.5\cdot7=1.5.$ (Because of the formula for the distance traveled in the $k$th one-second interval being $\Delta x_k = v_0 + a\left(k - \frac{1}{2}\right)$). The displacement at $t=30$ is 
 \[s=\frac{1}{2}at^2+v_0t=\frac{1}{2}\cdot7\cdot30^2+1.5\cdot30= \boxed{\textbf{(D)} ~3195}.\]~Bran\_Qin

 


	\item (\textit{AMC 10 10A 2021}) The quiz scores of a class with $k > 12$ students have a mean of $8$. The mean of a collection of $12$ of these quiz scores is $14$. What is the mean of the remaining quiz scores in terms of $k$?


 $\textbf{(A)} ~\frac{14-8}{k-12} \qquad\textbf{(B)} ~\frac{8k-168}{k-12} \qquad\textbf{(C)} ~\frac{14}{12} - \frac{8}{k} \qquad\textbf{(D)} ~\frac{14(k-12)}{k^2} \qquad\textbf{(E)} ~\frac{14(k-12)}{8k}$

 \textbf{Answer: }B



\textbf{Solution 1 (Generalized Value of k)}

The total score of the class is $8k,$ and the total score of the $12$ quizzes is $12\cdot14=168.$Therefore, for the remaining quizzes ($k-12$ of them), the total score is $8k-168.$ Their mean score is $\boxed{\textbf{(B)} ~\frac{8k-168}{k-12}}.$




\textbf{Solution 2 (Specified Values of k)}

Set $k=13.$ The answer is the same as the last student's quiz score, which is $8\cdot13-14\cdot12<0.$ From the answer choices, only $\boxed{\textbf{(B)} ~\frac{8k-168}{k-12}}$ is negative at $k=13.$




\textbf{Solution 3}

You know that the mean of the first $12$ students is $14,$ so that means all of them combined had a score of $12\cdot14 = 168.$ Set the mean of the remaining students (in other words the value you are trying to solve for), to $a.$ The total number of remaining students in a class of size $k$ can be written as $k-12.$ The total score $k-12$ students got combined can be written as $a(k-12),$ and the total score all of the students in the class got was $168 + a(k-12)$ (the first twelve students, plus the remaining students). The mean of the whole class can be written as $\frac{168 + a(k-12)}{k}.$ The mean of the class has already been given as $8,$ so by just writing the equation $\frac{168 + a(k-12)}{k} = 8,$ and solving for $a$ (the mean of $k-12$ students) will give you the answer in terms of $k,$ which is $\boxed{\textbf{(B)} ~\frac{8k-168}{k-12}}.$



	\item (\textit{AMC 10 10A 2021}) Chantal and Jean start hiking from a trailhead toward a fire tower. Jean is wearing a heavy backpack and walks slower. Chantal starts walking at $4$ miles per hour. Halfway to the tower, the trail becomes really steep, and Chantal slows down to $2$ miles per hour. After reaching the tower, she immediately turns around and descends the steep part of the trail at $3$ miles per hour. She meets Jean at the halfway point. What was Jean's average speed, in miles per hour, until they meet?


 $\textbf{(A)} ~\frac{12}{13} \qquad\textbf{(B)} ~1 \qquad\textbf{(C)} ~\frac{13}{12} \qquad\textbf{(D)}   ~\frac{24}{13} \qquad\textbf{(E)} ~2$

 \textbf{Answer: }A



\textbf{Solution 1 (Generalized Distance)}

Let $2d$ miles be the distance from the trailhead to the fire tower, where $d>0.$ When Chantal meets Jean, the two have traveled for 
 \[\frac d4 + \frac d2 + \frac d3 = d\left(\frac 14 + \frac 12 + \frac 13\right) =d\left(\frac{3}{12} + \frac{6}{12} + \frac{4}{12}\right)=\frac{13}{12}d\] hours. At that point, Jean has traveled for $d$ miles, so his average speed is $\frac{d}{\frac{13}{12}d}=\boxed{\textbf{(A)} ~\frac{12}{13}}$ miles per hour.

 



\textbf{Solution 2 (Specified Distance)}

<i>\textbf{We will follow the same template as shown in Solution 1, except that we will replace $\boldsymbol{d}$ with a convenient constant.}</i>

 Let $24$ miles be the distance from the trailhead to the fire tower. When Chantal meets Jean, the two have traveled for 
 \[\frac{12}{4} + \frac{12}{2}+\frac{12}{3}=3+6+4=13\] hours. At that point, Jean has traveled for $12$ miles, so his average speed is $\boxed{\textbf{(A)} ~\frac{12}{13}}$ miles per hour.

 



\textbf{Solution 3 (d=st)}

<i>\textbf{Shortened Solution}</i>
$d = st$, so Jean's speed can be represented as $s = \frac{d}{t}$. We know the time is $\left(\frac{d}{4} + \frac{d}{2} + \frac{d}{3}\right)$ (if the total distance is $2d$), so Jean's speed is $\frac{d}{(\frac{d}{4} + \frac{d}{2} + \frac{d}{3})} = \boxed{\textbf{(A)} ~\frac{12}{13}}$.

 <i>\textbf{Full Solution}</i>

 We know that distance traveled is equal to the speed multiplied with the time. So, $d=st$ and $t = \frac{d}{s}$. Let $2d$ be equal to the distance from the trailhead to the tower. Then, originally, Chantel travels a $d$ distance with at $4$ miles per hour. So, Chantel's time is $\frac{d}{4}$. From the midpoint to the tower, Chantel takes $\frac{d}{2}$ hours (since Chantel has speed of $2$ miles per hour.) Similarly, the time it takes for Chantel to return to the midpoint is $\frac{d}{3}$. Therefore, the total time is $\left(\frac{d}{4} + \frac{d}{2} + \frac{d}{3}\right)$. We can can substitute this time into the original equation of $d=st$ to obtain $d = s\left(\frac{d}{4} + \frac{d}{2} + \frac{d}{3}\right)$, so $\frac{d}{\left(\frac{d}{4} + \frac{d}{2} + \frac{d}{3}\right)} = s \implies \frac{1}{\frac{13}{12}} \implies \boxed{\textbf{(A)} ~\frac{12}{13}}$.

 


	\item (\textit{AMC 10 10A 2021}) Tom has a collection of $13$ snakes, $4$ of which are purple and $5$ of which are happy. He observes that


 
            \begin{itemize}
        \item all of his happy snakes can add,
            \end{itemize}
            
            \begin{itemize}
        \item none of his purple snakes can subtract, and
            \end{itemize}
            
            \begin{itemize}
        \item all of his snakes that can't subtract also can't add.
            \end{itemize}
            Which of these conclusions can be drawn about Tom's snakes?


 $\textbf{(A) }$ Purple snakes can add.


 $\textbf{(B) }$ Purple snakes are happy.


 $\textbf{(C) }$ Snakes that can add are purple.


 $\textbf{(D) }$ Happy snakes are not purple.


 $\textbf{(E) }$ Happy snakes can't subtract.


 
 \textbf{Answer: }D



\textbf{Solution 1 (Comprehensive Explanation of Logic)}

We are given that
 \begin{align*} \text{happy}&\Longrightarrow\text{can add}, &(1) \\ \text{purple}&\Longrightarrow\text{cannot subtract}, \hspace{15mm} &(2) \\ \text{cannot subtract}&\Longrightarrow\text{cannot add}. &(3) \end{align*}
Two solutions follow from here:

 



\textbf{Solution 1.1 (Intuitive)}

Combining $(2)$ and $(3)$ gives
 \begin{align*} \text{happy}&\Longrightarrow\text{can add}, &(1) \\ \lefteqn{\underbrace{\phantom{\text{purple}\Longrightarrow\text{cannot subtract}}}_{(2)}}\text{purple}&\Longrightarrow\overbrace{\text{cannot subtract}\Longrightarrow\text{cannot add}}^{(3)}. \hspace{2.5mm} &(*) \end{align*}
Clearly, the answer is $\boxed{\textbf{(D)}}.$




\textbf{Solution 1.2 (Rigorous)}

<i>\textbf{Recall that every conditional statement $\boldsymbol{p\Longrightarrow q}$ is always logically equivalent to its contrapositive $\boldsymbol{\lnot q\Longrightarrow\lnot p.}$}</i>

 Combining $(1),(2)$ and $(3)$ gives 
 \[\lefteqn{\underbrace{\phantom{\text{purple}\Longrightarrow\text{cannot subtract}}}_{(2)}}\text{purple}\Longrightarrow\overbrace{\text{cannot subtract}\Longrightarrow\lefteqn{\underbrace{\phantom{\text{cannot add}\Longrightarrow\text{not happy}}}_{\text{Contrapositive of }(1)}}\text{cannot add}}^{(3)}\Longrightarrow\text{not happy}. \hspace{15mm}(**)\] Applying the hypothetical syllogism to $(**),$ we conclude that 
 \[\text{purple}\Longrightarrow\text{not happy},\] whose contrapositive is 
 \[\text{happy}\Longrightarrow\text{not purple}.\] Therefore, the answer is $\boxed{\textbf{(D)}}.$
<u>\textbf{Remark}</u>

 The conclusions in the other choices do not follow from $(**):$
$\textbf{(A) }\text{purple}\Longrightarrow\text{can add}$
$\textbf{(B) }\text{purple}\Longrightarrow\text{happy}$
$\textbf{(C) }\text{can add}\Longrightarrow\text{purple}$
$\textbf{(E) }\text{happy}\Longrightarrow\text{cannot subtract}$




\textbf{Solution 2 (Process of Elimination)}

From Solution 1.1, we can also see this through the process of elimination. Statement $A$ is false because purple snakes cannot add. $B$ is false as well  because since happy snakes can add and purple snakes can not add, purple snakes are not happy snakes. $E$ is false using the same reasoning, purple snakes are not happy snakes so happy snakes can subtract since purple snakes cannot subtract. $C$ is false since snakes that can add are happy, not purple. That leaves statement D. $\boxed{\textbf{(D)}}$ is the only correct statement.

 



\textbf{Solution 3 (Rigorous)}

We first convert each statement to "If X, then Y" form:

 
            \begin{itemize}
        \item If a snake is happy, then it can add.
            \end{itemize}
            
            \begin{itemize}
        \item If a snake is purple, then it can't subtract.
            \end{itemize}
            
            \begin{itemize}
        \item If a snake can't subtract, then it can't add.
            \end{itemize}
            Now, we simply check the truth value for each statement: 

 
            \begin{enumerate} 
        \item Combining the last two propositions, we have
            \begin{itemize}<li>If a snake is purple, then it can't add. 

 </li>
            \end{itemize}
            Thus, $\textbf{(A)}$ is never true.

        \item From the last part, we found that
            \begin{itemize}<li>If a snake is purple, then it can't add. 

 </li>
            \end{itemize}
            Also, since the contrapositive of a proposition has the same truth value as the proposition itself, we know, from the first statement, that 
            \begin{itemize}<li>If a snake can't add, then it isn't happy. 

 </li>
            \end{itemize}
            Combining these two propositions, we find that 
            \begin{itemize}<li>If a snake is purple, then it isn't happy. Purple snakes are not happy.</li>

            \end{itemize}
            Thus, $\textbf{(B)}$ is never true.

        \item From part $\textbf{(A)},$ we found that "If a snake is purple, then it can't add." This implies its contrapositive, "If a snake can add, then it is not purple." is true, meaning $\textbf{(C)}$ is NEVER true. [Thanks again to MRENTHUSIASM for pointing this out!]

        \item From the first statement, we have 
            \begin{itemize}<li>If a snake is happy, then it can add.

 </li>
            \end{itemize}
            From the contrapositive of the third statement, we have
            \begin{itemize}<li>If a snake can add, then it can subtract. 

 </li>
            \end{itemize}
            Then, from the contrapositive of the second statement, we have 
            \begin{itemize}<li>If a snake can subtract, then it is not purple. 

 </li>
            \end{itemize}
            Combining all of these yields 
            \begin{itemize}<li>If a snake is happy, then it is not purple.

 </li>
            \end{itemize}
            Thus, $\textbf{(D)}$ is always true.

        \item From the first proposition, we have 
            \begin{itemize}<li>If a snake is happy, then it can add. 

 </li>
            \end{itemize}
            From the contrapositive of the third proposition, we have 
            \begin{itemize}<li>If a snake can add, then it can subtract. 

 </li>
            \end{itemize}
            Combining these two propositions gives 
            \begin{itemize}<li>If a snake is happy, then it can subtract. 

 </li>
            \end{itemize}
            Thus, $\textbf{(E)}$ is never true.

            \end{enumerate}
            Therefore, $\boxed{\textbf{(D)}}$ is our answer.

 


	\item (\textit{AMC 10 10A 2021}) When a student multiplied the number $66$ by the repeating decimal, 

 \[\underline{1}.\underline{a} \ \underline{b} \ \underline{a} \ \underline{b}\ldots=\underline{1}.\overline{\underline{a} \ \underline{b}},\] 
where $a$ and $b$ are digits, he did not notice the notation and just multiplied $66$ times $\underline{1}.\underline{a} \ \underline{b}.$ Later he found that his answer is $0.5$ less than the correct answer. What is the $2$-digit number $\underline{a} \ \underline{b}?$
$\textbf{(A) }15 \qquad \textbf{(B) }30 \qquad \textbf{(C) }45 \qquad \textbf{(D) }60 \qquad \textbf{(E) }75$

 \textbf{Answer: }E



\textbf{Solution 1}

We are given that $66\Bigl(\underline{1}.\overline{\underline{a} \ \underline{b}}\Bigr)-0.5=66\Bigl(\underline{1}.\underline{a} \ \underline{b}\Bigr),$ from which
 \begin{align*} 66\Bigl(\underline{1}.\overline{\underline{a} \ \underline{b}}\Bigr)-66\Bigl(\underline{1}.\underline{a} \ \underline{b}\Bigr)&=0.5 \\ 66\Bigl(\underline{1}.\overline{\underline{a} \ \underline{b}} - \underline{1}.\underline{a} \ \underline{b}\Bigr)&=0.5 \\ 66\Bigl(\underline{0}.\underline{0} \ \underline{0} \ \overline{\underline{a} \ \underline{b}}\Bigr)&=0.5 \\ 66\left(\frac{1}{100}\cdot\underline{0}.\overline{\underline{a} \ \underline{b}}\right)&=\frac12 \\ \underline{0}.\overline{\underline{a} \ \underline{b}}&=\frac{25}{33} \\ \underline{0}.\overline{\underline{a} \ \underline{b}}&=0.\overline{75} \\ \underline{a} \ \underline{b}&=\boxed{\textbf{(E) }75}. \end{align*}
~MRENTHUSIASM

 



\textbf{Solution 2}

It is known that $\underline{0}.\overline{\underline{a} \ \underline{b}}=\frac{\underline{a} \ \underline{b}}{99}$ and $\underline{0}.\underline{a} \ \underline{b}=\frac{\underline{a} \ \underline{b}}{100}.$ 

 Let $x=\underline{a} \ \underline{b}.$ We have 
 \[66\biggl(1+\frac{x}{99}\biggr)-66\biggl(1+\frac{x}{100}\biggr)=0.5.\] Expanding and simplifying give $\frac{x}{150}=0.5,$ so $x=\boxed{\textbf{(E) }75}.$




\textbf{Solution 3 (Similar to Solution 2)}

We have 
 \[66 \cdot \left(1 + \frac{10a+b}{100}\right) + \frac{1}{2} = 66 \cdot \left(1+ \frac{10a+b}{99}\right).\]Expanding both sides, we have 
 \[66 + \frac{33(10a+b)}{50} + \frac{1}{2} = 66 + \frac{2(10a+b)}{3}.\]Subtracting $66$ from both sides, we have 
 \[\frac{33(10a+b)}{50} + \frac{1}{2} = \frac{2(10a+b)}{3}.\]Multiplying both sides by $50 \cdot 3 = 150,$ we have 
 \[99(10a+b) + 75 = 100(10a+b).\]Thus, the answer is $10a+b = \boxed{\textbf{(E) }75}.$

 By letting $x=\underline{a} \ \underline{b}=10a+b,$ this solution is similar to Solution 2. In this solution, we solve for $10a+b$ as a whole.

 -mathboy282 (Solution)

 



\textbf{Solution 4 (Answer Choices \& Modular)}

Let $\underline{a} \ \underline{b}$ represent the two-digit number $ab$, not the product of the digits $a$ and $b$. We can construct fractions for the values $1.\underline{a} \ \underline{b}$, and $1.\overline{\underline{a} \ \underline{b}}$, which are $\frac{1\underline{a} \ \underline{b}}{100}$ and $\frac{1\underline{a} \ \underline{b}}{99}$  respectively. Multiplying by $66$ on both sides and adding $1/2$ to $\frac{1\underline{a} \ \underline{b}}{100}$ and simplifying, we get this: 
 \[\frac{33 * \underline{a}\underline{b} + 25}{50} = \frac{2\underline{a}\underline{b}}{3}.\]

 Looking at the answer choices, we notice that all of them are divisible by $3$. This means that since the right-hand side will result in an integer, the left-hand side needs to as well. This means that the numerator of the left-hand side fraction has to be divisible by $50$. So, we get this expression:

 $33 * \underline{a}\underline{b}\ + 25 \equiv 0\pmod{50}.$
$33 * \underline{a}\underline{b}\ \equiv -25\pmod{50}.$
$33 * \underline{a}\underline{b}\ \equiv 25\pmod{50}.$

 This means that the product of $33$ and $\underline{a}\underline{b}$ must have a remainder of $25$ when divided by $50$. Since it must have a remainder of $25$, the product should have a units digit of $5$, which eliminates $\textbf{(B) }$ and $\textbf{(D) }$. Multiplying $33$ to the rest of the answer choices, the only one which fills this requirement is $75$, which is $\boxed{\textbf{(E) }75}.$




\textbf{Solution 5 (A different way to calculate)}

$1.\overline{ab}=1+\dfrac{\overline{ab}_{10}}{10^2-1}=1+\dfrac{10a+b}{99}$ and $1.ab=1+\dfrac{a}{10}+\dfrac{b}{100}$, so our equation is 
 \[66\times\left(1+\dfrac{10a+b}{99}\right)=66\times\left(1+\dfrac{a}{10}+\dfrac{b}{100}\right)+0.5,\]so now let's simplify!\begin{align*}66\times\left(1+\dfrac{10a+b}{99}\right)&=66\times\left(1+\dfrac{a}{10}+\dfrac{b}{100}\right)+0.5\\1+\dfrac{10a+b}{99}&=1+\dfrac{a}{10}+\dfrac{b}{100}+\dfrac1{132}\\\dfrac{10a+b}{99}&=\dfrac{10a+b}{100}+\dfrac1{132}.\end{align*}
Notice that the $2$-digit number $\overline{ab}$ is just $10a+b$. Therefore, let's just find the value of $10a+b$.\newpage\begin{align*}\left(\dfrac1{99}-\dfrac1{100}\right)(10a+b)&=\dfrac1{132}\\10a+b&=\dfrac{\dfrac1{132}}{\dfrac1{99}-\dfrac1{100}}\\10a+b&=\dfrac{\dfrac1{132}}{\dfrac{100-99}{99\times100}}\\10a+b&=\dfrac{9900}{132}\\10a+b&=75.\end{align*}
We just got $10a+b=75$, and $75$ is indeed a $2$-digit number, so the two-digit number $\overline{ab}$ is $\boxed{75}$.

 


	\item (\textit{AMC 10 10A 2021}) What is the least possible value of $(xy-1)^2 + (x+y)^2$ for real numbers $x$ and $y$?


 $\textbf{(A)} ~0\qquad\textbf{(B)} ~\frac{1}{4}\qquad\textbf{(C)} ~\frac{1}{2}\qquad\textbf{(D)} ~1\qquad\textbf{(E)} ~2$

 \textbf{Answer: }D



\textbf{Solution 1 (Expand)}

Expanding, we get that the expression is $x^2+2xy+y^2+x^2y^2-2xy+1$ or $x^2+y^2+x^2y^2+1$. By the Trivial Inequality (all squares are nonnegative) the minimum value for this is $\boxed{\textbf{(D)} ~1}$, which can be achieved at $x=y=0$. 

 



\textbf{Solution 2 (Expand and then Factor)}

We expand the original expression, then factor the result by grouping:
 \begin{align*} (xy-1)^2+(x+y)^2&=\left(x^2y^2-2xy+1\right)+\left(x^2+2xy+y^2\right) \\ &=x^2y^2+x^2+y^2+1 \\ &=x^2\left(y^2+1\right)+\left(y^2+1\right) \\ &=\left(x^2+1\right)\left(y^2+1\right). \end{align*}
Clearly, both factors are positive. By the Trivial Inequality, we have 
 \[\left(x^2+1\right)\left(y^2+1\right)\geq\left(0+1\right)\left(0+1\right)=\boxed{\textbf{(D)} ~1}.\]Note that the least possible value of $(xy-1)^2+(x+y)^2$ occurs at $x=y=0.$




\textbf{Solution 3 (Beyond Overkill)}

Like solution 1, expand and simplify the original equation to $x^2+y^2+x^2y^2+1$ and let $f(x, y) = x^2+y^2+x^2y^2+1$. To find local extrema, find where $\nabla f(x, y) = \boldsymbol{0}$. First, find the first partial derivative with respect to x and y and find where they are $0$:
 \[\frac{\partial f}{\partial x} = 2x + 2xy^{2} = 2x(1 + y^{2}) = 0 \implies x = 0\]
 \[\frac{\partial f}{\partial y} = 2y + 2yx^{2} = 2y(1 + x^{2}) = 0 \implies y = 0\]Thus, there is a local extremum at $(0, 0)$. Because this is the only extremum, we can assume that this is a minimum because the problem asks for the minimum (though this can also be proven using the partial second derivative test) and the global minimum since it's the only minimum, meaning $f(0, 0)$ is the minimum of $f(x, y)$. Plugging $(0, 0)$ into $f(x, y)$, we find 1 $\implies \boxed{\textbf{(D)} ~1}$.

 Remark:Honestly, this isn't even that overkill. I did it this way, and it was not very bashy.~vaishnav

 


	\item (\textit{AMC 10 10A 2021}) Which of the following is equivalent to 
 \[(2+3)(2^2+3^2)(2^4+3^4)(2^8+3^8)(2^{16}+3^{16})(2^{32}+3^{32})(2^{64}+3^{64})?\]
$\textbf{(A)} ~3^{127} + 2^{127} \qquad\textbf{(B)} ~3^{127} + 2^{127} + 2 \cdot 3^{63} + 3 \cdot 2^{63} \qquad\textbf{(C)} ~3^{128}-2^{128} \qquad\textbf{(D)} ~3^{128} + 2^{128} \qquad\textbf{(E)} ~5^{127}$

 \textbf{Answer: }C



\textbf{Solution 1}

By multiplying the entire equation by $3-2=1$, all the terms will simplify by difference of squares, and the final answer is $\boxed{\textbf{(C)} ~3^{128}-2^{128}}$.

 Additionally, we could also multiply the entire equation (we can let it be equal to $x$) by $2-3=-1$. The terms again simplify by difference of squares. This time, we get $-x=2^{128}-3^{128} \Rightarrow x=3^{128}-2^{128}$. Both solutions yield the same answer.

 Note: Also notice when you multiply it by the first pair $(2+3)$, it immediately factors. Notice how it "domino" effects to the others, ultimately collapsing into $\boxed{\textbf{(C)} ~3^{128}-2^{128}}$.

 



\textbf{Solution 2}

If you weren't able to come up with the $(3 - 2)$ insight, then you could just notice that the answer is divisible by $2 + 3 = 5$, and $2^2 + 3^2 = 13$. We can then use Fermat's Little Theorem for $p = 5, 13$ on the answer choices to determine which of the answer choices are divisible by both $5$ and $13$. This is $\boxed{\textbf{(C)} ~3^{128}-2^{128}}$.

 -mewto

 



\textbf{Solution 3 (Units Digit and Multiples)}

Notice that simplifying the first term gives 5, which means that the units digit of the entire expression must be 5. Also, note how the numbers appear to form multiples of 2, so therefore our answer should contain a power of 2. The only one that satisfies the above properties is $\boxed{\textbf{(C)} ~3^{128}-2^{128}}$.

 



\textbf{Solution 4}

After expanding the first few terms, the result after each term appears to be $2^{2^n-1} + 2^{2^n-2}\cdot{3^1} + 2^{2^n-3}\cdot{3^2} + \cdots + 2^1\cdot{3^{2^n-2}} + 3^{2^n-1}$ where $n$ is the number of terms expanded. We can prove this using mathematical induction. The base step is trivial. When expanding another term, all of the previous terms multiplied by $2^{2^{n-1}}$ would give $2^{2^n-1} + 2^{2^n-2}\cdot{3^1} + 2^{2^n-3}\cdot{3^2} + \cdots + 2^{2^{n-1}+1}\cdot{3^{2^{n-1}-1}} + 2^{2^{n-1}}\cdot{3^{2^{n-1}}}$, and all the previous terms multiplied by $3^{2^{n-1}}$ would give $3^{2^n-1} + 3^{2^n-2}\cdot{2^1} + 3^{2^n-3}\cdot{2^2} + \cdots + 3^{2^{n-1}+1}\cdot{2^{2^{n-1}-1}} + 3^{2^{n-1}}\cdot{2^{2^{n-1}}}$. Their sum is equal to $2^{2^n-1} + 2^{2^n-2}\cdot{3^1} + 2^{2^n-3}\cdot{3^2} + \cdots + 2^1\cdot{3^{2^n-2}} + 3^{2^n-1}$, so the proof is complete. Since $\frac{3^{2^n}-2^{2^n}}{3-2}$ is equal to $2^{2^n-1} + 2^{2^n-2}\cdot{3^1} + 2^{2^n-3}\cdot{3^2} + \cdots + 2^1\cdot{3^{2^n-2}} + 3^{2^n-1}$, the answer is $\frac{3^{2^7}-2^{2^7}}{3-2}=\boxed{\textbf{(C)} ~3^{128}-2^{128}}$.

 



\textbf{Solution 5 (Engineer's Induction)}

We can compute some of the first few partial products, and notice that $\prod_{k = 0}^{2^n} (2^{2^n}+3^{2^n}) = 3^{2^{n+1}} - 2^{2^{n+1}}$. As we don't have to prove this, we get the product is $3^{2^7} - 2^{2^7} = 3^{128} - 2^{128}$, and click $\boxed{\textbf{(C)} ~3^{128}-2^{128}}$. 

 



\textbf{Solution 6 (Difference of Squares)}

We notice that the first term is equal to $3^2 - 2^2$. If we multiply this by the second term, then we will get $(3^2 - 2^2)(3^2 + 2^2)$, and we can simplify by using difference of squares to obtain $3^4 - 2^4$. If we multiply this by the third term and simplify using difference of squares again, we get $3^8 - 2^8$. We can continue down the line until we multiply by the last term, $3^{64} + 2^{64}$, and get $\boxed{\textbf{(C)} ~3^{128}-2^{128}}$.

 



\textbf{Solution 7 (Generalization)}

Let’s generalize to $x=3$ and $y=2$. Then we get:
 \[(y+x)(y^2+x^2)(y^4+x^4)(y^8+x^8)(y^{16}+x^{16})(y^{32}+x^{32})(y^{64}+x^{64}).\]We see that the first term is $y+x$, and the next is $y^2+x^2$. We realize that if we multiply the first term by $y-x$, the result by difference of squares will be $y^2-x^2$: we can proceed to use difference of squares on that. In other words, the equation has the domino effect, and all we need to get started is to multiply the whole equation by $y-x$ (keeping in mind to divide $y-x$ at the end.)

 Now we get:
 \begin{align*} &\hspace{5.125mm}(y-x)(y+x)(y^2+x^2)(y^4+x^4)(y^8+x^8)(y^{16}+x^{16})(y^{32}+x^{32})(y^{64}+x^{64}) \\ &=(y^2-x^2)(y^2+x^2)(y^4+x^4)(y^8+x^8)(y^{16}+x^{16})(y^{32}+x^{32})(y^{64}+x^{64}) \\ &=(y^4-x^4)(y^4+x^4)(y^8+x^8)(y^{16}+x^{16})(y^{32}+x^{32})(y^{64}+x^{64}) \\ & \ \vdots \\ &=(y^{64}-x^{64})(y^{64}+x^{64}) \\ &=y^{128}-x^{128}. \end{align*}
Now we can plug in $y = 2$ and $x = 3:$ 
 \[(2-3)(2+3)(2^2+3^2)(2^4+3^4)(2^8+3^8)(2^{16}+3^{16})(2^{32}+3^{32})(2^{64}+3^{64}) = 2^{128}-3^{128}.\]However, we must not forget to divide by $2-3 = -1$ at the end! Dividing, we get the answer of 
 \[3^{128}-2^{128} = \boxed{\textbf{(C)} ~3^{128}-2^{128}}.\]~KingRavi

 



\textbf{Solution 8 (Generalization)}

More generally, we have 
 \[\bigl(a^n-b^n\bigr)\bigl(a^n+b^n\bigr)=a^{2n}-b^{2n}\hspace{20mm}(\bigstar)\] by the difference of squares.

 Multiplying the original expression by $1=3-2$ and then applying $(\bigstar)$ repeatedly, we get
 \begin{align*} &\hspace{5.125mm}\bigl(2+3\bigr)\bigl(2^2+3^2\bigr)\bigl(2^4+3^4\bigr)\bigl(2^8+3^8\bigr)\bigl(2^{16}+3^{16}\bigr)\bigl(2^{32}+3^{32}\bigr)\bigl(2^{64}+3^{64}\bigr) \\ &=\underbrace{\bigl(\hspace{3.25mm}1\hspace{3.25mm}\bigr)}_{3-2}\underbrace{\bigl(2+3\bigr)\bigl(2^2+3^2\bigr)\bigl(2^4+3^4\bigr)\bigl(2^8+3^8\bigr)\bigl(2^{16}+3^{16}\bigr)\bigl(2^{32}+3^{32}\bigr)\bigl(2^{64}+3^{64}\bigr)}_{\text{Apply the }\textbf{Commutative Property of Addition}\text{ to all factors.}} \\ &=\underbrace{\bigl(3-2\bigr)\bigl(3+2\bigr)}_{\substack{\text{Apply }(\bigstar) \\ \text{to get }3^2-2^2.}}\bigl(3^2+2^2\bigr)\bigl(3^4+2^4\bigr)\bigl(3^8+2^8\bigr)\bigl(3^{16}+2^{16}\bigr)\bigl(3^{32}+2^{32}\bigr)\bigl(3^{64}+2^{64}\bigr) \\ &=\underbrace{\bigl(3^2-2^2\bigr)\bigl(3^2+2^2\bigr)}_{\substack{\text{Apply }(\bigstar) \\ \text{to get }3^4-2^4.}}\bigl(3^4+2^4\bigr)\bigl(3^8+2^8\bigr)\bigl(3^{16}+2^{16}\bigr)\bigl(3^{32}+2^{32}\bigr)\bigl(3^{64}+2^{64}\bigr) \\ &=\underbrace{\bigl(3^4-2^4\bigr)\bigl(3^4+2^4\bigr)}_{\substack{\text{Apply }(\bigstar) \\ \text{to get }3^8-2^8.}}\bigl(3^8+2^8\bigr)\bigl(3^{16}+2^{16}\bigr)\bigl(3^{32}+2^{32}\bigr)\bigl(3^{64}+2^{64}\bigr) \\ & \ \vdots \\ &=\boxed{\textbf{(C)} ~3^{128}-2^{128}}. \end{align*}
~MRENTHUSIASM

 



\textbf{Solution 9 (Desperate & Cheese, Similar to Solution 5)}

One would probably assume that the factorization that this is equal to can be generalized for the pattern 
 \[(2+3)(2^2 + 3^2)(2^4+3^4)\cdots(2^n + 3^n).\]So we can take a smaller portion and apply the answer choices there. We get
 \[(2+3)(2^2 + 3^2) = (5)(13) = 65.\]We notice that the patterns have a relationship where its either the sum of the powers or the next power not included that is incorporated. So we change the numbers from $127$ to $3$, $128$ to $4$, and $63$ to $1$ (since $63$ is the sum of the powers, divided by $2$, rounded down). Checking the values of the other patterns, the only one that gives $65$ is $3^4 - 2^4$, which means the answer is $\boxed{\textbf{(C)} ~3^{128}-2^{128}}.$



	\item (\textit{AMC 10 10A 2021}) For which of the following integers $b$ is the base-$b$ number $2021_b - 221_b$ not divisible by $3$?


 $\textbf{(A)} ~3 \qquad\textbf{(B)} ~4\qquad\textbf{(C)} ~6\qquad\textbf{(D)} ~7\qquad\textbf{(E)} ~8$

 \textbf{Answer: }E



\textbf{Solution 1 (Factor)}

We have 
 \begin{align*} 2021_b - 221_b &= (2021_b - 21_b) - (221_b - 21_b) \\ &= 2000_b - 200_b \\ &= 2b^3 - 2b^2 \\ &= 2b^2(b-1), \end{align*}
which is divisible by $3$ \textbf{\textit{unless}} $b\equiv2\pmod{3}.$ The only choice congruent to $2$ modulo $3$ is $\boxed{\textbf{(E)} ~8}.$




\textbf{Solution 2 (Vertical Subtraction)}

Vertically subtracting $2021_b - 221_b,$ we see that the ones place becomes $0,$ and so does the $b^1$ place. Then, we perform a carry (make sure the carry is in base $b$). Let $b-2 = A.$ Then, we have our final number as 
 \[1A00_b.\]Now, when expanding, we see that this number is simply $b^3 - (b - 2)^2.$

 Now, notice that the final number will only be congruent to 
 \[b^3-(b-2)^2\equiv0\pmod{3}.\] If either $b\equiv0\pmod{3},$ or if $b\equiv1\pmod{3}$ (because note that $(b - 2)^2$ would become $\equiv1\pmod{3},$ and $b^3$ would become $\equiv1\pmod{3}$ as well, and therefore the final expression would become $1-1\equiv0\pmod{3}.$ Therefore, $b$ must be $\equiv2\pmod{3}.$ Among the answers, only $8$ is $\equiv2\pmod{3},$ and therefore our answer is $\boxed{\textbf{(E)} ~8}.$




\textbf{Solution 3 (Answer Choices)}

By the definition of bases, we have 
 \[2021_b - 221_b = \left(2b^3+2b+1\right) - \left(2b^2+2b+1\right).\]For values $b_1$ and $b_2$ such that $b_1\equiv b_2\pmod{3},$ we get 
 \[\left(2b_1^3+2b_1+1\right) - \left(2b_1^2+2b_1+1\right) \equiv \left(2b_2^3+2b_2+1\right) - \left(2b_2^2+2b_2+1\right) \pmod{3}.\]Note that answer choices $\textbf{(A)},\textbf{(B)},\textbf{(C)},\textbf{(D)},\textbf{(E)}$ are congruent to $0,1,0,1,2$ modulo $3,$ respectively. So, $\textbf{(A)}$ and $\textbf{(C)}$ are either both correct or both incorrect. Since there is only one correct answer, $\textbf{(A)}$ and $\textbf{(C)}$ are both incorrect. Similarly, $\textbf{(B)}$ and $\textbf{(D)}$ are both incorrect. This leaves us with $\boxed{\textbf{(E)} ~8},$ the answer choice with a unique residue modulo $3.$




\textbf{Solution 4 (like Solution 1)}

As with the beginning of Solution 1, we get that 3 does not divide $2b^2(b-1)$. We can just test each answer choice (doing short circuit evaluation for $3$ and $6$), and finding that only $\boxed{\textbf{(E)} ~8}$ works. Honestly, Solution 1 is easier, but during an actual test when your brain is fried, I'd just check each option.

 


 



\textbf{Solution 5}

Test all the answer choices and see that $8 \textbf{ (E)}$ is the answer!

 


	\item (\textit{AMC 10 10A 2021}) Two right circular cones with vertices facing down as shown in the figure below contain the same amount of liquid. The radii of the tops of the liquid surfaces are $3 \text{ cm}$ and $6 \text{ cm}$. Into each cone is dropped a spherical marble of radius $1 \text{ cm}$, which sinks to the bottom and is completely submerged without spilling any liquid. What is the ratio of the rise of the liquid level in the narrow cone to the rise of the liquid level in the wide cone? 


 
                    \begin{figure}[H]
                        \centering
                        \includegraphics[width=0.5\textwidth]{Problem/AMC_10_10A_2021_12_Problem_Figure_1.png}
                    \end{figure}
                    
$\textbf{(A) }1:1 \qquad \textbf{(B) }47:43 \qquad \textbf{(C) }2:1 \qquad \textbf{(D) }40:13 \qquad \textbf{(E) }4:1$

 \textbf{Answer: }E



\textbf{Solution 1 (Algebra)}

\textbf{<u>Initial Scenario</u>}

 Let the heights of the narrow cone and the wide cone be $h_1$ and $h_2,$ respectively. We have the following table:
 \[\begin{array}{cccccc} & \textbf{Base Radius} & \textbf{Height} & & \textbf{Volume} & \\ [2ex] \textbf{Narrow Cone} & 3 & h_1 & & \frac13\pi(3)^2h_1=3\pi h_1 & \\ [2ex] \textbf{Wide Cone} & 6 & h_2 & & \hspace{2mm}\frac13\pi(6)^2h_2=12\pi h_2 & \end{array}\]Equating the volumes gives $3\pi h_1=12\pi h_2,$ which simplifies to $\frac{h_1}{h_2}=4.$ 

 Furthermore, by similar triangles:

 
            \begin{itemize}
        \item For the narrow cone, the ratio of the base radius to the height is $\frac{3}{h_1},$ which always remains constant.
            \end{itemize}
            
            \begin{itemize}
        \item For the wide cone, the ratio of the base radius to the height is $\frac{6}{h_2},$ which always remains constant.
            \end{itemize}
            Two solutions follow from here:

 



\textbf{Solution 1.1 (Properties of Fractions)}

\textbf{<u>Final Scenario</u>}

 For the narrow cone and the wide cone, let their base radii be $3x$ and $6y$ (for some $x,y>1$), respectively. By the similar triangles discussed above, their heights must be $h_1x$ and $h_2y,$ respectively. We have the following table:
 \[\begin{array}{cccccc} & \textbf{Base Radius} & \textbf{Height} & & \textbf{Volume} & \\ [2ex] \textbf{Narrow Cone} & 3x & h_1x & & \frac13\pi(3x)^2(h_1x)=3\pi h_1 x^3 & \\ [2ex] \textbf{Wide Cone} & 6y & h_2y & & \hspace{2.0625mm}\frac13\pi(6y)^2(h_2y)=12\pi h_2 y^3 & \end{array}\]Recall that $\frac{h_1}{h_2}=4.$ Equating the volumes gives $3\pi h_1 x^3=12\pi h_2 y^3,$ which simplifies to $x^3=y^3,$ or $x=y.$

 Finally, the requested ratio is 
 \[\frac{h_1 x - h_1}{h_2 y - h_2}=\frac{h_1 (x-1)}{h_2 (y-1)}=\frac{h_1}{h_2}=\boxed{\textbf{(E) }4:1}.\]<u>\textbf{Remarks}</u>

            \begin{enumerate} 
        \item This solution uses the following property of fractions:  For unequal positive numbers $a,b,c$ and $d,$ if $\frac ab = \frac cd = k,$ then $\frac{a\pm c}{b\pm d}=\frac{bk\pm dk}{b\pm d}=\frac{(b\pm d)k}{b\pm d}=k.$


        \item This solution shows that, regardless of the shape or the volume of the solid dropped into each cone, the requested ratio stays the same as long as the solid sinks to the bottom and is completely submerged without spilling any liquid.

            \end{enumerate}
            



\textbf{Solution 1.2 (Bash)}

\textbf{<u>Final Scenario</u>}

 For the narrow cone and the wide cone, let their base radii be $r_1$ and $r_2,$ respectively; let their rises of the liquid levels be $\Delta h_1$ and $\Delta h_2,$ respectively. We have the following table:
 \[\begin{array}{cccccc} & \textbf{Base Radius} & \textbf{Height} & & \textbf{Volume} & \\ [2ex]  \textbf{Narrow Cone} & r_1 & h_1+\Delta h_1 & & \frac13\pi r_1^2(h_1+\Delta h_1) & \\ [2ex]  \textbf{Wide Cone} & r_2 & h_2+\Delta h_2 & & \frac13\pi r_2^2(h_2+\Delta h_2) & \end{array}\]By the similar triangles discussed above, we get
 \begin{align*} \frac{3}{h_1}&=\frac{r_1}{h_1+\Delta h_1} &\implies \quad r_1&=\frac{3}{h_1}(h_1+\Delta h_1), & \hspace{10mm} (1) \\ \frac{6}{h_2}&=\frac{r_2}{h_2+\Delta h_2} &\implies \quad r_2&=\frac{6}{h_2}(h_2+\Delta h_2). & (2) \end{align*}
The volume of the marble dropped into each cone is $\frac43\pi(1)^3=\frac43\pi.$

 Now, we set up an equation for the volume of the narrow cone, then express $\Delta h_1$ in terms of $h_1:$
 \begin{align*} \frac13\pi r_1^2(h_1+\Delta h_1) &= 3\pi h_1+\frac43\pi \\ \frac13 r_1^2(h_1+\Delta h_1) &= 3h_1+\frac43 \\ \frac13\left(\frac{3}{h_1}(h_1+\Delta h_1)\right)^2(h_1+\Delta h_1) &= 3h_1+\frac43 &&\text{by }(1) \\ \frac{3}{h_1^2}(h_1+\Delta h_1)^3 &= 3h_1+\frac43 \\ (h_1+\Delta h_1)^3 &= h_1^3 + \frac{4h_1^2}{9} \\ \Delta h_1 &= \sqrt[3]{h_1^3 + \frac{4h_1^2}{9}}-h_1. \end{align*}
Next, we set up an equation for the volume of the wide cone, then express $\Delta h_2$ in terms of $h_2:$
 \[\frac13\pi r_2^2(h_2+\Delta h_2) = 12\pi h_2+\frac43\pi.\]Using a similar process from above, we get 
 \[\Delta h_2 = \sqrt[3]{h_2^3+\frac{h_2^2}{9}}-h_2.\]Recall that $\frac{h_1}{h_2}=4.$ Therefore, the requested ratio is 
 \begin{align*} \frac{\Delta h_1}{\Delta h_2}&=\frac{\sqrt[3]{h_1^3 + \frac{4h_1^2}{9}}-h_1}{\sqrt[3]{h_2^3+\frac{h_2^2}{9}}-h_2} \\ &=\frac{\sqrt[3]{(4h_2)^3 + \frac{4(4h_2)^2}{9}}-4h_2}{\sqrt[3]{h_2^3+\frac{h_2^2}{9}}-h_2} \\ &=\frac{\sqrt[3]{4^3\left(h_2^3 + \frac{h_2^2}{9}\right)}-4h_2}{\sqrt[3]{h_2^3+\frac{h_2^2}{9}}-h_2} \\ &=\frac{4\sqrt[3]{h_2^3+\frac{h_2^2}{9}}-4h_2}{\sqrt[3]{h_2^3+\frac{h_2^2}{9}}-h_2} \\ &=\boxed{\textbf{(E) }4:1}. \end{align*}
~MRENTHUSIASM

 



\textbf{Solution 2 (Approximate Cones with Cylinders)}

The heights of the cones are not given, so suppose the heights are very large (i.e. tending towards infinity) in order to approximate the cones as cylinders with base radii $3$ and $6$ and infinitely large height. Then the base area of the wide cylinder is $4$ times that of the narrow cylinder. Since we are dropping a ball of the same volume into each cylinder, the water level in the narrow cone/cylinder should rise $\boxed{\textbf{(E) } 4}$ times as much.

 



\textbf{Solution 3 (Calculus)}

The volume of the shorter cone can be expressed as $V_1=3\pi h_1$, and the volume of the larger cone can be expressed as $V_2=12\pi h_2$. We also know that the volume changes by $\frac{4\pi}{3}$, because the volume of the $1$cm sphere is $\frac{4\pi}{3}$.

 Taking the derivative of the equations, we get: $dV_1=3\pi(dh_1)=\frac{4\pi}{3}$ and $dV_2=12\pi(dh_2)=\frac{4\pi}{3}$.

 Therefore, $dh_1=\frac{4\pi}{3}\cdot\frac{1}{3\pi}=\frac{4}{9}$ and $dh_2=\frac{4\pi}{3}\cdot\frac{1}{12\pi}=\frac{1}{9}$. The ratio is $\frac{4}{9}:\frac{1}{9}$ giving us the answer of $\boxed{\textbf{(E) }4:1}$.

 



\textbf{Solution 4 (Extremely Quick Observation)}

Note that as long as the volumes are equal, the ratio of the heights of Cone 2 to Cone 1 is always $1:4$. We can prove this as follows:

 If we fix Cone 2 to have a certain height, then the volume of the cone is $\frac{1}{3}\cdot36\pi{h_1}=12\pi{h_1}$. Now if Cone 1 has the same volume, its height would be $\frac{12\pi{h_1}}{\frac{1}{3}\cdot9\pi}=\frac{12\pi{h_1}}{3\pi}=4h_1$.

 So the answer is $\boxed{\textbf{(E) } 4 : 1}$

 Minor edits by: \href{https://artofproblemsolving.com/wiki/index.php?title=User:T3chn0b14d3&action=edit&redlink=1}{T3CHN0B14D3}

 \textbf{Remark} You can cheese this by letting $r_1$ and $r_2$ be some small integers. I recommend $4$ and $1.$



	\item (\textit{AMC 10 10A 2021}) What is the volume of tetrahedron $ABCD$ with edge lengths $AB = 2$, $AC = 3$, $AD = 4$, $BC = \sqrt{13}$, $BD = 2\sqrt{5}$, and $CD = 5$ ?


 $\textbf{(A)} ~3 \qquad\textbf{(B)} ~2\sqrt{3} \qquad\textbf{(C)} ~4\qquad\textbf{(D)} ~3\sqrt{3}\qquad\textbf{(E)} ~6$

 \textbf{Answer: }C



\textbf{Solution 1 (Three Right Triangles)}

Drawing the tetrahedron out and testing side lengths, we realize that the $\triangle ACD, \triangle ABC,$ and $\triangle ABD$ are right triangles by the Converse of the Pythagorean Theorem. It is now easy to calculate the volume of the tetrahedron using the formula for the volume of a pyramid. If we take $\triangle ADC$ as the base, then $\overline{AB}$ must be the altitude. The volume of tetrahedron $ABCD$ is $\dfrac{1}{3} \cdot \dfrac{3 \cdot 4}{2} \cdot 2=\boxed{\textbf{(C)} ~4}.$This solution does not have a reasoning for using triangle ADC as base~Icewolf10 ~Bakedpotato66 ~MRENTHUSIASM

 



\textbf{Solution 2 (One Right Triangle)}

We will place tetrahedron $ABCD$ in the $xyz$-plane. By the Converse of the Pythagorean Theorem, we know that $\triangle ACD$ is a right triangle. Without the loss of generality, let $A=(0,0,0), C=(3,0,0), D=(0,4,0),$ and $B=(x,y,z).$

 We apply the Distance Formula to $\overline{BA},\overline{BC},$ and $\overline{BD},$ respectively:
 \begin{align*} x^2+y^2+z^2&=2^2, &(1) \\ (x-3)^2+y^2+z^2&=\sqrt{13}^2, &(2) \\ x^2+(y-4)^2+z^2&=\left(2\sqrt5\right)^2. &\hspace{1mm} (3) \end{align*}
Subtracting $(1)$ from $(2)$ gives $-6x+9=9,$ from which $x=0.$

 Subtracting $(1)$ from $(3)$ gives $-8y+16=16,$ from which $y=0.$

 Substituting $(x,y)=(0,0)$ into $(1)$ produces $z^2=4,$ or $|z|=2.$

 Let the brackets denote areas. Finally, we find the volume of tetrahedron $ABCD$ using $\triangle ACD$ as the base:
 \begin{align*} V_{ABCD}&=\frac13\cdot[ACD]\cdot h_B \\ &=\frac13\cdot\left(\frac12\cdot AC \cdot AD\right)\cdot |z| \\ &=\boxed{\textbf{(C)} ~4}. \end{align*}
~MRENTHUSIASM

 



\textbf{Solution 3 (Trirectangular Tetrahedron)}

\href{https://mathworld.wolfram.com/TrirectangularTetrahedron.html}{https://mathworld.wolfram.com/TrirectangularTetrahedron.html}

 Given the observations from Solution 1, where $\triangle ACD, \triangle ABC,$ and $\triangle ABD$ are right triangles, the base is $\triangle ABD.$ We can apply the information about a trirectangular tetrahedron (all of the face angles are right angles), which states that the volume is
 \begin{align*} V&=\frac16\cdot AB\cdot AC\cdot BD \\ &=\frac16\cdot2\cdot4\cdot3 \\ &=\boxed{\textbf{(C)} ~4}. \end{align*}
~AMC60 (Solution)

 \textbf{Remark} Here is a similar problem from another AMC test: \href{https://artofproblemsolving.com/wiki/index.php?title=2015\_AMC\_10A\_Problems/Problem\_21}{2015 AMC 10A Problem 21}.

 


	\item (\textit{AMC 10 10A 2021}) All the roots of the polynomial $z^6-10z^5+Az^4+Bz^3+Cz^2+Dz+16$ are positive integers, possibly repeated. What is the value of $B$?


 $\textbf{(A) }{-}88 \qquad \textbf{(B) }{-}80 \qquad \textbf{(C) }{-}64 \qquad \textbf{(D) }{-}41\qquad \textbf{(E) }{-}40$

 \textbf{Answer: }A



\textbf{Solution 1}

By Vieta's formulas, the sum of the six roots is $10$ and the product of the six roots is $16$. By inspection, we see the roots are $1, 1, 2, 2, 2,$ and $2$, so the function is $(z-1)^2(z-2)^4=(z^2-2z+1)(z^4-8z^3+24z^2-32z+16)$. Therefore, calculating just the $z^3$ terms, we get $B = -32 - 48 - 8 = \boxed{\textbf{(A) }{-}88}$.

 



\textbf{Solution 2}

Using the same method as Solution 1, we find that the roots are $2, 2, 2, 2, 1,$ and $1$. Note that $B$ is the negation of the 3rd symmetric sum of the roots. Using casework on the number of 1's in each of the $\binom {6}{3} = 20$ products $r_a \cdot r_b \cdot r_c,$ we obtain 
 \[B= - \left(\binom {4}{3} \binom {2}{0} \cdot 2^{3} + \binom {4}{2} \binom{2}{1} \cdot 2^{2} \cdot 1 + \binom {4}{1} \binom {2}{2} \cdot 2 \right) = -\left(32+48+8 \right) = \boxed{\textbf{(A) }{-}88}.\]~ike.chen

 



\textbf{Solution 3 (brute force)}

Same as solution 1, we find the roots are $1, 1, 2, 2, 2, 2$ by Vieta's formula. This next part is brute force expansion.$-B=\\r_1r_2r_3+r_1r_2r_4+r_1r_2r_5+r_1r_2r_6+\\r_1r_3r_4+r_1r_3r_5+r_1r_3r_6+\\r_1r_4r_5+r_1r_4r_6+\\r_1r_5r_6+\\r_2r_3r_4+r_2r_3r_5+r_2r_3r_6+\\r_2r_4r_5+r_2r_4r_6+\\r_2r_5r_6+\\r_3r_4r_5+r_3r_4r_6+\\r_3r_5r_6+\\r_4r_5r_6\\=2 \times 4 + 4 \times 12 + 8 \times 4\\ =8 + 48 + 32\\B = \boxed{\textbf{(A) }{-}88}$



	\item (\textit{AMC 10 10A 2021}) Values for $A,B,C,$ and $D$ are to be selected from $\{1, 2, 3, 4, 5, 6\}$ without replacement (i.e. no two letters have the same value). How many ways are there to make such choices so that the two curves $y=Ax^2+B$ and $y=Cx^2+D$ intersect? (The order in which the curves are listed does not matter; for example, the choices $A=3, B=2, C=4, D=1$ is considered the same as the choices $A=4, B=1, C=3, D=2.$)


 $\textbf{(A) }30 \qquad \textbf{(B) }60 \qquad \textbf{(C) }90 \qquad \textbf{(D) }180 \qquad \textbf{(E) }360$

 \textbf{Answer: }C



\textbf{Solution 1 (Intuition):}

Visualizing the two curves, we realize they are both parabolas with the same axis of symmetry. WLOG the first equation is above the second, since order doesn't matter. Then $C>A$ and $B>D$. Therefore the number of ways to choose the four integers is $\tbinom{6}{2}\tbinom{4}{2}=90$, and the answer is $\boxed{\textbf{(C) }90}$. 

 



\textbf{Solution 2 (Algebra):}

Setting $y = Ax^2+B = Cx^2+D$, we find that $Ax^2-Cx^2 = x^2(A-C) = D-B$, so $x^2 = \frac {D-B}{A-C} \ge 0$ by the trivial inequality. This implies that $D-B$ and $A-C$ must both be positive or negative. If two distinct values are chosen for $(A, C)$ and $(B, D)$ respectively, there are $2$ ways to order them so that both the numerator and denominator are positive/negative (increasing and decreasing). We must divide by $2$ at the end, however, since the $2$ curves aren't considered distinct. Calculating, we get 
 \[\frac {1}{2} \cdot \binom {6}{2} \binom {4}{2} \cdot 2 = \boxed{\textbf{(C) }90}.\] ~ike.chen

 



\textbf{Solution 3 (Symmetry):}

Like in \href{https://artofproblemsolving.com\#Solution\_2\_(Algebra):}{Solution 2}, we find $\frac {D-B}{A-C} \ge 0$. Notice that, since $D \ne B$, this expression can never equal $0$, and since $A \ne C$, there won't be a divide-by-$0$. This means that every choice results in either a positive or a negative value.

 For every choice of $(A,B,C,D)$ that results in a positive value, we can flip $B$ and $D$ to obtain a corresponding negative value. This is a bijection (we could flip $B$ and $D$ again to obtain the original choice (injectivity) and we could flip $B$ and $D$ from any negative choice to obtain the corresponding positive choice (surjectivity)), so half of the choices are positive (where the curves intersect) and half are negative (where they don't).

 This means that of the $\frac{6\cdot5\cdot4\cdot3}{2} = 180$ total choices (dividing by $2$ because the order of the curves does not matter), half of them, or $\frac{180}{2} = \boxed{\textbf{(C) }90}$, lead to intersecting curves.

 



\textbf{Solution 4 (Simple)}

Notice that one of the parabolas must be the wider one and the other one must be the thinner one. There are only two options: the wider one is above the thinner one, or the thinner one is above the wider one. Only the first option works. Therefore, out of the $\frac{1}{2}\cdot 6\cdot 5\cdot 4\cdot 3=180$ ways to pick $A, B, C,$ and $D$ without regard to order, only $\frac{1}{2}\cdot 180=90$ of these work.

 



\textbf{Solution 5}

We notice that the only way this will work is if the wider parabola is lower and the thinner one is higher. We have $15$ ways to pick values for one parabola to be thinner than the other; we multiply this by $2$ to get $30$ because there are $2$ choices for which parabola is thinner. After picking these $2$ numbers, there are $3$ ways to choose values for the thinner parabola to have a higher origin than the wider one. We multiply $30$ by $3$ to get the answer of $\boxed{\textbf{(C) }90}$.

 


	\item (\textit{AMC 10 10A 2021}) In the following list of numbers, the integer $n$ appears $n$ times in the list for $1 \leq n \leq 200$.
 \[1, 2, 2, 3, 3, 3, 4, 4, 4, 4, \dots, 200, 200, \dots , 200\]What is the median of the numbers in this list?


 $\textbf{(A)} ~100.5\qquad\textbf{(B)} ~134\qquad\textbf{(C)} ~142\qquad\textbf{(D)} ~150.5\qquad\textbf{(E)} ~167$

 \textbf{Answer: }C



\textbf{Solution 1}

There are $1+2+..+199+200=\frac{(200)(201)}{2}=20100$ numbers in total. Let the median be $k$. We want to find the median $k$ such that
 \[\frac{k(k+1)}{2}=20100/2,\]or
 \[k(k+1)=20100.\]Note that $\sqrt{20100} \approx 142$. Plugging this value in as $k$ gives
 \[\frac{1}{2}(142)(143)=10153.\]$10153-142<10050$, so $142$ is the $152$nd and $153$rd numbers, and hence, our desired answer. $\fbox{(C) 142}$.

 Note that we can derive $\sqrt{20100} \approx 142$ through the formula 
 \[\sqrt{n} = \sqrt{a+b} \approx \sqrt{a} + \frac{b}{2\sqrt{a} + 1},\]where $a$ is a perfect square less than or equal to $n$. We set $a$ to $19600$, so $\sqrt{a} = 140$, and $b = 500$. We then have $n \approx 140 + \frac{500}{2(140)+1} \approx 142$. ~approximation by ciceronii

 Note by Fasolinka (use answer choices): Once you know that the answer is in the 140s range (200,000 is around 14\textasciicircum{}2 times 10\textasciicircum{}2) by the approximation, it is highly improbable for the answer to be anything but C.

 



\textbf{Solution 2}

The $x$th number of this sequence is $\left\lceil\frac{-1\pm\sqrt{1+8x}}{2}\right\rceil$ via the quadratic formula. We can see that if we halve $x$ we end up getting $\left\lceil\frac{-1\pm\sqrt{1+4x}}{2}\right\rceil$. This is approximately the number divided by $\sqrt{2}$. $\frac{200}{\sqrt{2}} = 141.4$ and since $142$ looks like the only number close to it, it is answer $\boxed{(C) 142}$ ~Lopkiloinm

 



\textbf{Solution 3 (Answer Choices)}

We can look at answer choice $\textbf{(C)}$, which is $142$ first. That means that the number of numbers from $1$ to $142$ is roughly the number of numbers from $143$ to $200$.

 The number of numbers from $1$ to $142$ is $\frac{142(142+1)}{2}$ which is approximately $10000.$ The number of numbers from $143$ to $200$ is $\frac{200(200+1)}{2}-\frac{142(142+1)}{2}$ which is approximately $10000$ as well. Therefore, we can be relatively sure the answer choice is $\boxed{\textbf{(C)} ~142}.$

 -\href{https://artofproblemsolving.com/wiki/index.php?title=User:PureSwag}{PureSwag}

 



\textbf{Solution 4 (Geometry)}

We can arrange the numbers in the following pattern:
 \[ \begin{array}{cccccc} \ &\ &\ &\ &\ 200 & \\ \ &\ &\ &\ 199 & \ 200 & \\ \ &\ &\ \iddots& \ \vdots& \ \vdots& \\ \ &\ 2& \ \cdots& \ 199& \ 200& \\ 1 & \ 2 & \ \cdots& \ 199& \ 200& \end{array} \]

 We can see this as a isosceles right triangle, with legs of length $200.$ 
                    \begin{figure}[H]
                        \centering
                        \includegraphics[width=0.5\textwidth]{Solution/AMC_10_10A_2021_16_Solution_Figure_1.png}
                    \end{figure}
                    

 Let $x$ be the side length such that both sides of the triangle have the same area. The desired answer is then around $x$ because about half of the numbers in the list fall on each side.

 Solving for $x$ yields:
 \begin{align*} \frac{x^2}{2} =& \:\frac{1}{2} \cdot \frac{200^2}{2} \\ x^2 =& \:\frac{1}{2}\cdot 200^2 \\ x =& \:\frac{200}{\sqrt{2}} = \: 100\sqrt{2} \approx 141. \end{align*}
We see that $\boxed{(C) \: 142}$ is the closest to $x$ by far, and thus, can be relatively certain this is the answer. 

 


	\item (\textit{AMC 10 10A 2021}) Trapezoid $ABCD$ has $\overline{AB} \parallel \overline{CD}$, $BC = CD = 43$, and $\overline{AD} \perp \overline{BD}$. Let $O$ be the intersection of the diagonals $\overline{AC}$ and $\overline{BD}$, and let $P$ be the midpoint of $\overline{BD}$. Given that $OP = 11$, the length $AD$ can be written in the form $m\sqrt{n}$, where $m$ and $n$ are positive integers and $n$ is not divisible by the square of any prime. What is $m + n$?


 $\textbf{(A)} ~65\qquad\textbf{(B)} ~132\qquad\textbf{(C)} ~157\qquad\textbf{(D)} ~194\qquad\textbf{(E)} ~215$

 \textbf{Answer: }D



\textbf{Solution 1 (Similar Triangles and Pythagorean Theorem)}

Angle chasing* reveals that $\triangle BPC\sim\triangle BDA$, therefore 
 \[2=\frac{BD}{BP}=\frac{AB}{BC}=\frac{AB}{43},\]or $AB=86$.

 Additional angle chasing shows that $\triangle ABO\sim\triangle CDO$, therefore
 \[2=\frac{AB}{CD}=\frac{BO}{OD}=\frac{BP+11}{BP-11},\]or $BP=33$ and $BD=66$.

 Since $\triangle ADB$ is right, the Pythagorean theorem implies that
 \[AD=\sqrt{86^2-66^2}=4\sqrt{190}.\]The answer is $4+190=\boxed{\textbf{(D) }194}$.

 
            \begin{itemize}
        \item Angle Chasing: If we set $\angle DBC = \alpha$, then we know that $\angle DCB = 180^\circ-2\alpha$ because $\triangle DBC$ is isosceles. Then, $\angle BCP = 90^\circ-\alpha$ (since $P$ is the midpoint of $\overline{BD},$ and therefore line $\overline{CP}$ is a median, which then means $\angle BCP$ bisects $\angle DBC$), so $\angle BPC$ is a right angle. Because $\angle BDC = \alpha$ and $\overline{AB}\parallel\overline{DC}$, we conclude that $\angle ABD = \alpha$ too. Lastly, because $\triangle BPC$ and $\triangle BDA$ are both right triangles, they are similar by AA.
            \end{itemize}
            \textbf{=Another way to angle chase} It is easy to show that triangles $CDO$ and $ODA$ are similar through $AA$. From there, we already know that through isosceles triangles $CPD$ and $CPB$ are congruent. Hence we can simply see that angle $PCO$ and angle $OAD$ are congruent through similar triangles. Also, angle $OAB$ and angle $OCD$ are congruent through parallelism. Thus, angle $BAD$ and angle $DCP$ are congruent. Then, through $AA$, $\triangle BCP \cong \triangle DCP \sim \triangle BDA$.

 This skips the unnecessary step of angle chasing via similar angles of an isosceles triangle and skips the step of manually counting each angle's measure in terms of $\alpha$. However, no discredit to the original angle chasing comment, it has its merits as well (using the angle measures in the above angle chasing technique proves useful for showing other statements [and cooler facts] about the problem).

 



\textbf{Solution 2 (Similar Triangles, Areas, Pythagorean Theorem)}

Since $\triangle BCD$ is isosceles with base $\overline{BD},$ it follows that median $\overline{CP}$ is also an altitude. Let $OD=x$ and $CP=h,$ so $PB=x+11.$

 Since $\angle AOD=\angle COP$ by vertical angles, we conclude that $\triangle AOD\sim\triangle COP$ by AA, from which $\frac{AD}{CP}=\frac{OD}{OP},$ or 
 \[AD=CP\cdot\frac{OD}{OP}=h\cdot\frac{x}{11}.\] Let the brackets denote areas. Notice that $[AOD]=[BOC]$ (By the same base and height, we deduce that $[ACD]=[BDC].$ Subtracting $[OCD]$ from both sides gives $[AOD]=[BOC].$). Doubling both sides produces
 \begin{align*} 2[AOD]&=2[BOC] \\ OD\cdot AD&=OB\cdot CP \\ x\left(\frac{hx}{11}\right)&=(x+22)h \\ x^2&=11(x+22). \end{align*}
Rearranging and factoring result in $(x-22)(x+11)=0,$ from which $x=22.$

 Applying the Pythagorean Theorem to right $\triangle CPB,$ we have 
 \[h=\sqrt{43^2-33^2}=\sqrt{(43+33)(43-33)}=\sqrt{760}=2\sqrt{190}.\] Finally, we get 
 \[AD=h\cdot\frac{x}{11}=4\sqrt{190},\] so the answer is $4+190=\boxed{\textbf{(D) }194}.$




\textbf{Solution 3 (Short)}


                    \begin{figure}[H]
                        \centering
                        \includegraphics[width=0.5\textwidth]{Solution/AMC_10_10A_2021_17_Solution_Figure_1.png}
                    \end{figure}
                    

 


 Let $CP = y$. $CP$ a is perpendicular bisector of $DB.$ Then, let $DO = x,$ thus $DP = PB = 11+x.$

 (1) $\triangle CPO \sim \triangle ADO,$ so we get $\frac{AD}{x} = \frac{y}{11},$ or $AD = \frac{xy}{11}.$

 (2) Applying Pythagorean Theorem on $\triangle CDP$ gives $(11+x)^2 + y^2 = 43^2.$

 (3) $\triangle BPC \sim \triangle BDA$ with ratio $1:2,$ so $AD = 2y$ using the fact that $P$ is the midpoint of $BD$.

 Thus, $\frac{xy}{11} = 2y,$ or $x = 22.$ And $y = \sqrt{43^2 - 33^2} = 2 \sqrt{190},$ so $AD = 4 \sqrt{190}$ and the answer is $4+190=\boxed{\textbf{(D) }194}.$




\textbf{Solution 4 (Law of Sines and Similar Triangles)}

This solution will refer to the diagram section at the top of this page.

 
We know $\angle ABD = \angle BDC = \angle DBC$. Let $PB = 11 + x$. Notice that $\triangle ABD$ is similar to $\triangle DOC$, so $\frac{AB}{43} = \frac{22 + x}{x} \implies AB = \frac{43(22 + x)}{x}$. Now we apply the Law of Sines to isosceles triangle $\triangle DCB$:

 
 \[\frac{\sin(\angle ABD)}{43} = \frac{\sin(180^\circ - 2\angle ABD)}{22 + 2x} = \frac{2\sin(\angle ABD)\cos(\angle ABD)}{22 + 2x}\]

 \[\frac{1}{43} = \frac{\cos(\angle ABD)}{11 + x} \implies \cos(\angle ABD) = \frac{11 + x}{43}\]

 From right triangle $\triangle BAD$ we already know that $\cos(\angle ABD) = \frac{BD}{AB} = \frac{(22 + 2x)x}{43(22 + x)}$.

 Setting these equations equal to each other gives us:
 \[\frac{(22 + 2x)x}{43(22 + x)} = \frac{11 + x}{43}\]

 \[x = 22 \implies BD = 2(11 + x) = 66, AB = \frac{43 \cdot 44}{22} = 86.\]

 So $AD = \sqrt{86^2 - 66^2} = 4\sqrt{190}$. Ans = $\boxed{194}$.

 



\textbf{Solution 5 (Extending the Line)}

Observe that $\triangle BPC$ is congruent to $\triangle DPC$; both are similar to $\triangle BDA$. Let's extend $\overline{AD}$ and $\overline{BC}$ past points $D$ and $C$ respectively, such that they intersect at a point $E$. Observe that $\angle BDE$ is $90$ degrees, and that $\angle DBE \cong \angle PBC \cong \angle DBA \implies \angle DBE \cong \angle DBA$. Thus, by ASA, we know that $\triangle ABD \cong \triangle EBD$, thus, $AD = ED$, meaning $D$ is the midpoint of $AE$.Let $M$ be the midpoint of $\overline{DE}$. Note that $\triangle CME$ is congruent to $\triangle BPC$, thus $BC = CE$, meaning $C$ is the midpoint of $\overline{BE}.$

 Therefore, $\overline{AC}$ and $\overline{BD}$ are both medians of $\triangle ABE$. This means that $O$ is the centroid of $\triangle ABE$; therefore, because the centroid divides the median in a 2:1 ratio, $\frac{BO}{2} = DO = \frac{BD}{3}$. Recall that $P$ is the midpoint of $BD$; $DP = \frac{BD}{2}$. The question tells us that $OP = 11$; $DP-DO=11$; we can write this in terms of $DB$; $\frac{DB}{2}-\frac{DB}{3} = \frac{DB}{6} = 11 \implies DB = 66$.

 We are almost finished. Each side length of $\triangle ABD$ is twice as long as the corresponding side length $\triangle CBP$ or $\triangle CPD$, since those triangles are similar; this means that $AB = 2 \cdot 43 = 86$. Now, by Pythagorean theorem on $\triangle ABD$, $AB^{2} - BD^{2} = AD^{2} \implies 86^{2}-66^{2} = AD^{2} \implies AD = \sqrt{3040} \implies AD = 4 \sqrt{190}$. 

 The answer is $4+190 = \boxed{\textbf{(D) }194}$.

 



\textbf{Solution 6}

Since $P$ is the midpoint of isosceles triangle $BCD$, it would be pretty easy to see that $CP\perp BD$. Since $AD\perp BD$ as well, $AD\parallel CP$. Connecting $AP$, it’s obvious that $[ADC]=[ADP]$. Since $DP=BP$, $[APB]=[ADC]$. 

 Since $P$ is the midpoint of $BD$, the height of $\triangle APB$ on side $AB$ is half that of $\triangle ADC$ on $CD$. Since $[APB]=[ADC]$, $AB=2CD$. 

 As a basic property of a trapezoid, $\triangle AOB \sim \triangle COD$, so $\frac{OB}{OD}=\frac{AB}{CD}=2$, or $OB=2OD$. Letting $OD=x$, then $PB=DP=11+x$, and $OB=22+x$. Hence $22+x=2x$ and $x=22$.

 Since $\triangle AOD \sim \triangle COP$, $\frac{AD}{PC}=\frac{OD}{OP}=2$. Since $PD=11+22=33$, $PC=\sqrt{43^2-33^2}=\sqrt{760}$.

 So, $AD=2\sqrt{760}=4\sqrt{190}$. The correct answer is $\boxed{\textbf{(D) }194}$.

 



\textbf{Solution 7 (Coordinate Geometry)}

Let $D$ be the origin of the cartesian coordinate plane, $B$ lie on the positive $x$-axis, and $A$ lie on the negative $y$-axis. Then let the coordinates of $B = (2a,0), A = (0, -2b).$ Then the slope of $AB$ is $\frac{b}{a}.$ Since $AB \parallel CD$ the slope of $CD$ is the same. Note that as $\triangle DCB$ is isosceles $C$ lies on $x = a.$ Thus since $CD$ has equation $y = \frac{b}{a}x$ ($D$ is the origin), $C = (a,b).$ Therefore $AC$ has equation $y = \frac{3b}{a}x - 2b$ and intersects $BD$ ($x$-axis) at $O =\left(\frac{2}{3}a, 0\right).$ The midpoint of $BD$ is $P = (a,0),$ so $OP = \frac{a}{3} = 11,$ from which $a = 33.$ Then by Pythagorean theorem on $\triangle DPC$ ($\triangle DBC$ is isosceles), we have $b = \sqrt{43^2 - 33^2} = 2\sqrt{190},$ so $2b=4\sqrt{190}.$

 Finally, the answer is $4+190=\boxed{\textbf{(D) }194}.$




\textbf{Solution 8 (Trigonometry)}


                    \begin{figure}[H]
                        \centering
                        \includegraphics[width=0.5\textwidth]{Solution/AMC_10_10A_2021_17_Solution_Figure_2.png}
                    \end{figure}
                    set $\angle BDC  = \theta$
$BD = 2*DP = 2*43*\cos{\theta}$
$AB = \frac{BD}{\cos{\angle {DBA}}} = \frac{BD}{\cos{\theta}}=   2 \cdot 43 \cdot \frac{\cos{\theta}}{\cos{\theta}}=86$
$\frac{OP}{DO}=\frac{CP}{AD}$
$\frac{11}{43\cdot \cos{\theta}-11}=\frac{43\cdot\sin{\theta}}{86\cdot \sin{\theta}}$
$\cos{\theta}=\frac{33}{43}$
$AD=86\cdot \sin{\theta}=2 \sqrt{760}=4\sqrt{190}=\boxed{\textbf{(D)}194}$.

 


	\item (\textit{AMC 10 10A 2021}) Let $f$ be a function defined on the set of positive rational numbers with the property that $f(a\cdot b)=f(a)+f(b)$ for all positive rational numbers $a$ and $b$. Furthermore, suppose that $f$ also has the property that $f(p)=p$ for every prime number $p$. For which of the following numbers $x$ is $f(x)<0$?


 $\textbf{(A)} ~\frac{17}{32}\qquad\textbf{(B)} ~\frac{11}{16}\qquad\textbf{(C)} ~\frac{7}{9}\qquad\textbf{(D)} ~\frac{7}{6} \qquad\textbf{(E)} ~\frac{25}{11}$

 \textbf{Answer: }E



\textbf{Solution 1 (Intuitive)}

From the answer choices, note that
 \begin{align*} f(25)&=f\left(\frac{25}{11}\cdot11\right) \\ &=f\left(\frac{25}{11}\right)+f(11) \\ &=f\left(\frac{25}{11}\right)+11. \end{align*}
On the other hand, we have
 \begin{align*} f(25)&=f(5\cdot5) \\ &=f(5)+f(5) \\ &=5+5 \\ &=10. \end{align*}
Equating the expressions for $f(25)$ produces 
 \[f\left(\frac{25}{11}\right)+11=10,\] from which $f\left(\frac{25}{11}\right)=-1.$ Therefore, the answer is $\boxed{\textbf{(E) }\frac{25}{11}}.$
<u>\textbf{Remark}</u>

 Similarly, we can find the outputs of $f$ at the inputs of the other answer choices:
 \begin{alignat*}{10} &\textbf{(A)} \qquad && f\left(\frac{17}{32}\right) \quad && = \quad && 7 \\  &\textbf{(B)} \qquad && f\left(\frac{11}{16}\right) \quad && = \quad && 3 \\  &\textbf{(C)} \qquad && f\left(\frac{7}{9}\right) \quad && = \quad && 1 \\  &\textbf{(D)} \qquad && f\left(\frac{7}{6}\right) \quad && = \quad && 2 \end{alignat*}Alternatively, refer to Solutions 2 and 4 for the full processes.

 



\textbf{Solution 2 (Specific)}

We know that $f(p) = f(p \cdot 1) = f(p) + f(1)$. By transitive, we have 
 \[f(p) = f(p) + f(1).\]Subtracting $f(p)$ from both sides gives $0 = f(1).$Also
 \[f(2)+f\left(\frac{1}{2}\right)=f(1)=0 \implies 2+f\left(\frac{1}{2}\right)=0 \implies f\left(\frac{1}{2}\right) = -2\]
 \[f(3)+f\left(\frac{1}{3}\right)=f(1)=0 \implies 3+f\left(\frac{1}{3}\right)=0 \implies f\left(\frac{1}{3}\right) = -3\]
 \[f(11)+f\left(\frac{1}{11}\right)=f(1)=0 \implies 11+f\left(\frac{1}{11}\right)=0 \implies f\left(\frac{1}{11}\right) = -11\]In $\textbf{(A)}$ we have $f\left(\frac{17}{32}\right)=17+5f\left(\frac{1}{2}\right)=17-5(2)=7$.

 In $\textbf{(B)}$ we have $f\left(\frac{11}{16}\right)=11+4f\left(\frac{1}{2}\right)=11-4(2)=3$.

 In $\textbf{(C)}$ we have $f\left(\frac{7}{9}\right)=7+2f\left(\frac{1}{3}\right)=7-2(3)=1$.

 In $\textbf{(D)}$ we have $f\left(\frac{7}{6}\right)=7+f\left(\frac{1}{2}\right)+f\left(\frac{1}{3}\right)=7-2-3=2$.

 In $\textbf{(E)}$ we have $f\left(\frac{25}{11}\right)=10+f\left(\frac{1}{11}\right)=10-11=-1$.

 Thus, our answer is $\boxed{\textbf{(E) }\frac{25}{11}}$.

 



\textbf{Solution 3 (Generalized)}

Consider the rational $\frac{a}{b}$, for $a,b$ integers. We have $f(a)=f\left(\frac{a}{b}\cdot b\right)=f\left(\frac{a}{b}\right)+f(b)$. So $f\left(\frac{a}{b}\right)=f(a)-f(b)$. Let $p$ be a prime. Notice that $f(p^k)=kf(p)$. And $f(p)=p$. So if $a=p_1^{a_1}p_2^{a_2}\cdots p_k^{a_k}$, $f(a)=a_1p_1+a_2p_2+\cdots+a_kp_k$. We simply need this to be greater than what we have for $f(b)$. Notice that for answer choices $\textbf{(A)},\textbf{(B)},\textbf{(C)},$ and $\textbf{(D)}$, the numerator has fewer prime factors than the denominator, and so they are less likely to work. We check $\textbf{(E)}$ first, and it works, therefore the answer is $\boxed{\textbf{(E) }\frac{25}{11}}$.  

 



\textbf{Solution 4 (Generalized)}

We derive the following properties of $f:$

            \begin{enumerate} 
        \item By induction, we have 
 \[f\left(\prod_{k=1}^{n}a_k\right)=\sum_{k=1}^{n}f(a_k)\] for all positive rational numbers $a_k$ and positive integers $n.$ Since positive powers are just repeated multiplication of the base, it follows that 
 \[f\left(a^n\right)=f\left(\prod_{k=1}^{n}a\right)=\sum_{k=1}^{n}f(a)=nf(a)\] for all positive rational numbers $a$ and positive integers $n.$


        \item For all positive rational numbers $a,$ we have 
 \[f(a)=f(a\cdot1)=f(a)+f(1),\] from which $f(1)=0.$

        \item For all positive rational numbers $a,$ we have 
 \[f(a)+f\left(\frac1a\right)=f\left(a\cdot\frac1a\right)=f(1)=0,\] from which $f\left({\frac 1a}\right)=-f(a).$

            \end{enumerate}
            For all positive integers $x$ and $y,$ suppose $\prod_{k=1}^{m}p_k^{d_k}$ and $\prod_{k=1}^{n}q_k^{e_k}$ are their respective prime factorizations. We get
 \begin{align*} f\left(\frac xy\right)&=f(x)+f\left(\frac 1y\right) \\ &=f(x)-f(y) && \hspace{10mm}\text{by Property 3} \\ &=f\left(\prod_{k=1}^{m}p_k^{d_k}\right)-f\left(\prod_{k=1}^{n}q_k^{e_k}\right) \\ &=\left[\sum_{k=1}^{m}f\left(p_k^{d_k}\right)\right]-\left[\sum_{k=1}^{n}f\left(q_k^{e_k}\right)\right] && \hspace{10mm}\text{by Property 1} \\ &=\left[\sum_{k=1}^{m}d_k f\left(p_k\right)\right]-\left[\sum_{k=1}^{n}e_k f\left(q_k\right)\right] && \hspace{10mm}\text{by Property 1} \\ &=\left[\sum_{k=1}^{m}d_k p_k \right]-\left[\sum_{k=1}^{n}e_k q_k \right]. \end{align*}
We apply $f$ to each fraction in the answer choices:
 \begin{alignat*}{10} &\textbf{(A)} \qquad && f\left(\frac{17}{32}\right) \quad && = \quad && f\left(\frac{17\textasciicircum{}1}{2\textasciicircum{}5}\right) \quad && = \quad && [1(17)]-[5(2)] \quad && = \quad && 7 \\  &\textbf{(B)} \qquad && f\left(\frac{11}{16}\right) \quad && = \quad && f\left(\frac{11\textasciicircum{}1}{2\textasciicircum{}4}\right) \quad && = \quad && [1(11)]-[4(2)] \quad && = \quad && 3 \\  &\textbf{(C)} \qquad && f\left(\frac{7}{9}\right) \quad && = \quad && f\left(\frac{7\textasciicircum{}1}{3\textasciicircum{}2}\right)  \quad && = \quad && [1(7)]-[2(3)]  \quad && = \quad && 1 \\  &\textbf{(D)} \qquad && f\left(\frac{7}{6}\right) \quad && = \quad && f\left(\frac{7\textasciicircum{}1}{2\textasciicircum{}1\cdot3\textasciicircum{}1}\right) \quad && = \quad && [1(7)]-[1(2)+1(3)] \quad && = \quad && 2 \\ &\textbf{(E)} \qquad && f\left(\frac{25}{11}\right) \quad && = \quad && f\left(\frac{5\textasciicircum{}2}{11\textasciicircum{}1}\right) \quad && = \quad && [2(5)]-[1(11)] \quad && = \quad && {-}1 \end{alignat*}Therefore, the answer is $\boxed{\textbf{(E) }\frac{25}{11}}.$




\textbf{Solution 5 (Quick, Dirty, and Frantic Last Hope)}

Note that answer choices $\textbf{(A)}$ through $\textbf{(D)}$ are $\frac{\text{prime}}{\text{composite}},$ whereas $\textbf{(E)}$ is $\frac{\text{composite}}{\text{prime}}.$ Because the functional equation is related to primes, we hope that the uniqueness of answer choice $\boxed{\textbf{(E) }\frac{25}{11}}$ is enough.

 



\textbf{Solution 6 (Rushed Generalization)}

If f(a $\cdot$ b) = f(a) + f(b), and if f(p) = p, then f(p $\cdot$ p) = 2p. You can do this multiple times (Ex: f(p\textasciicircum{}3) = 3p). You can quickly assume then, that f(p\textasciicircum{}n) = np. Thus the answer choices can then be rewritten as the product of a prime and another prime to the negative power. Answer choices A-C are straightforward. For D, you can rewrite $\frac{1}{6}$ as $\frac{1}{2}$ $\cdot$ $\frac{1}{3}$. When you get to E, you get f(25) + f($\frac{1}{11}$), which is 10 - 11, which is -1. So the answer is $\boxed{\textbf{(E) }\frac{25}{11}}$




\textbf{Solution 7 (Generalized)}

Note that for each of the answer choices we can multiply the fractions by their denominators to be left with only the numerator and also have the prime factorization of the denominators, then set the function equal to the numerator. This will be shown throughout the solution.

 For answer choice $\text{(A) }\dfrac{17}{32}$, we have $f\left(\dfrac{17}{32}\right)$. Now, we can set up an equation by multiplying $\dfrac{17}{32}$ by $32$ and setting it equal to $17$. This will give $f\left(32\cdot\dfrac{17}{32}\right)=f(17)\rightarrow f(32)+f\left(\dfrac{17}{32}\right)=f(17)$. Our goal is to find $f\left(\dfrac{17}{32}\right)$, therefore we must find $f(32)$ and $f(17)$. Since $f(p)=p$ for any prime $p$, $f(17)=17$. Taking the prime factorization of $32$ gives $2^5$, so $f(32)=f(2\cdot2\cdot2\cdot2\cdot2)=f(2)+f(2)+f(2)+f(2)+f(2)=5f(2)=5\cdot2=10$. Therefore, $10+f\left(\dfrac{17}{32}\right)=17$ and $f\left(\dfrac{17}{32}\right)=7$
$\text{(B) }\dfrac{11}{16}$: $f\left(\dfrac{11}{16}\right)\rightarrow f\left(16\cdot\dfrac{11}{16}\right)=f(11)$$\rightarrow f(16)+f\left(\dfrac{11}{16}\right)=11\rightarrow 4f(2)+f\left(\dfrac{11}{16}\right)=11$$\rightarrow f\left(\dfrac{11}{16}\right)=3$
$\text{(C) }\dfrac{7}{9}$: $f\left(\dfrac{7}{9}\right)\rightarrow f\left(9\cdot\dfrac{7}{9}\right)=f(7)$$\rightarrow f(9)+f\left(\dfrac{7}{9}\right)=7\rightarrow 2f(3)+f\left(\dfrac{7}{9}\right)=7$$\rightarrow f\left(\dfrac{7}{9}\right)=1$
$\text{(D) }\dfrac{7}{6}$: $f\left(\dfrac{7}{6}\right)\rightarrow f\left(6\cdot\dfrac{7}{6}\right)=f(7)$$\rightarrow f(6)+f\left(\dfrac{7}{6}\right)=7\rightarrow f(2)+f(3)+f\left(\dfrac{7}{6}\right)=7$$\rightarrow f\left(\dfrac{7}{6}\right)=2$
$\text{(E) }\dfrac{25}{11}$: $f\left(\dfrac{25}{11}\right)\rightarrow f\left(11\cdot\dfrac{25}{11}\right)=f(25)$$\rightarrow f(11)+f\left(\dfrac{25}{11}\right)=2f(5)\rightarrow 11+f\left(\dfrac{25}{11}\right)=10$$\rightarrow f\left(\dfrac{25}{11}\right)=-1$

 Therefore, the answer is $\boxed{\text{(E) }\dfrac{25}{11}}$

 Note: The general strategy here was the setting up of equations to find $f(x)$. By setting it equal to $f(a)$ where $a$ was an integer, we could take the prime factorization to find the value of $f(a)$ and also set up an equation involving $f(\text{denominator})$ because the denominator was also an integer, therefore we could take the prime factorization and find its value

 



\textbf{Solution 8 (Generalized But Fast)}

Let $p$ be a prime number. We have that $f(p)=p.$ Now, consider $f\left(\frac1{p^a}\right)$ where $a$ is a positive integer. Notice that
 \[f\left(p^{a+1}\cdot\frac1{p^a}\right)=f(p)=p\quad\text{and}\quad f\left(p^{a+1}\cdot\frac1{p^a}\right)=f(p^{a+1})+f\left(\frac1{p^a}\right)=(a+1)p+f\left(\frac1{p^a}\right)\]and hence
 \[(a+1)p+f\left(\frac1{p^a}\right)=p\implies f\left(\frac1{p^a}\right)=-ap.\]Thus, we see that
 \[f\left(\frac{25}{11}\right)=2f(5)+f\left(\frac1{11}\right)=10-11=-1<0\implies \boxed{\textbf{(E)}\,\frac{25}{11}}\]~shockfront99

 


	\item (\textit{AMC 10 10A 2021}) The area of the region bounded by the graph of 
 \[x^2+y^2 = 3|x-y| + 3|x+y|\]is $m+n\pi$, where $m$ and $n$ are integers. What is $m + n$?


 $\textbf{(A)} ~18\qquad\textbf{(B)} ~27\qquad\textbf{(C)} ~36\qquad\textbf{(D)} ~45\qquad\textbf{(E)} ~54$

 \textbf{Answer: }E



\textbf{Solution 1}

In order to attack this problem, we can use casework on the sign of $|x-y|$ and $|x+y|$.

 Case 1: $|x-y|=x-y, |x+y|=x+y$

 Substituting and simplifying, we have $x^2-6x+y^2=0$, i.e. $(x-3)^2+y^2=3^2$, which gives us a circle of radius $3$ centered at $(3,0)$.

 Case 2: $|x-y|=y-x, |x+y|=x+y$

 Substituting and simplifying again, we have $x^2+y^2-6y=0$, i.e. $x^2+(y-3)^2=3^2$. This gives us a circle of radius $3$ centered at $(0,3)$.

 Case 3: $|x-y|=x-y, |x+y|=-x-y$

 Doing the same process as before, we have $x^2+y^2+6y=0$, i.e. $x^2+(y+3)^2=3^2$. This gives us a circle of radius $3$ centered at $(0,-3)$.

 Case 4: $|x-y|=y-x, |x+y|=-x-y$

 One last time: we have $x^2+y^2+6x=0$, i.e. $(x+3)^2+y^2=3^2$. This gives us a circle of radius $3$ centered at $(-3,0)$.

 After combining all the cases and drawing them on the Cartesian Plane, this is what the diagram looks like:

 
                    \begin{figure}[H]
                        \centering
                        \includegraphics[width=0.5\textwidth]{Solution/AMC_10_10A_2021_19_Solution_Figure_1.png}
                    \end{figure}
                    Now, the area of the shaded region is just a square with side length $6$ with four semicircles of radius $3$.The area is $6\cdot6+4\cdot \frac{9\pi}{2} = 36+18\pi$. The answer is $36+18$ which is $\boxed{\textbf{(E) }54}$




\textbf{Solution 2}

A somewhat faster variant of solution 1 is to use a bit of symmetry in order to show that the remaining three cases are identical to Case 1 in the above solution, up to rotations by $90^{\circ}$ about the origin. This allows us to quickly sketch the region after solving Case 1.

 Upon simplifying Case 1, we obtain $(x-3)^2 + y^2 = 3^2$ which is a circle of radius 3 centered at $(3,0)$. We remark that only the points on the semicircle where $x \ge 3$ work here, since Case 1 assumes $x-y \ge 0$ and $x+y \ge 0$. Next, we observe that an ordered pair is a solution to the given equation if and only if any of its $90^{\circ}$ rotations about the origin is a solution. This follows as the value of $x^2+y^2-3(|x-y|+|x+y|)$ is invariant to $90^{\circ}$ rotations, since $x^2+y^2$ simply represents the square of the distance to the origin (which is unchanged upon rotation), and $|x-y|+|x+y|$ is the sum of the distances to the lines $y=x$ and $y=-x$, multiplied by $\sqrt{2}$ (also unchanged upon $90^{\circ}$ rotation).

 By the above observation, we can quickly sketch the remainder of the region, and the area is $\boxed{\textbf{(E) }54}$ as above.

 



\textbf{Solution 3 (Guessing)}

Assume $y$ = $0$. We get that $x$ = $6$. That means that this figure must contain the points $(0,6), (6,0), (0, -6), (-6, 0)$. Now, assume that $x$ = $y$. We get that $x$ = $3 \sqrt 3$. We get the points $(3,3), (3,-3), (-3, 3), (-3, -3)$. 

 Since this contains $x^2 + y^2$, assume that there are circles. Therefore, we can guess that there is a center square with area $6 \cdot 6$ = $36$ and $4$ semicircles with radius $3$. We get $4$ semicircles with area $4.5 \pi$, and therefore the answer is $36+18$ = $\boxed {(E)54}$

 \textbf{Remark} This problem asks for the area of \textbf{the union of these four circles}:

 
                    \begin{figure}[H]
                        \centering
                        \includegraphics[width=0.5\textwidth]{Solution/AMC_10_10A_2021_19_Solution_Figure_2.png}
                    \end{figure}
                    \textbf{Rotation by 45 degrees} with the help of rotation \href{https://www.wolframalpha.com/input/?i=rotate+45+degrees}{https://www.wolframalpha.com/input/?i=rotate+45+degrees}, we can simplify the equation to $a^2 + b^2 = 3\sqrt2|a| + 3\sqrt2|b|$. Then follow the previous question\href{https://artofproblemsolving.com/wiki/index.php/2016\_AMC\_10B\_Problems/Problem\_21}{https://artofproblemsolving.com/wiki/index.php/2016\_AMC\_10B\_Problems/Problem\_21}

 


	\item (\textit{AMC 10 10A 2021}) In how many ways can the sequence $1, 2, 3, 4, 5$ be rearranged so that no three consecutive terms are increasing and no three consecutive terms are decreasing?


 $\textbf{(A)} ~10\qquad\textbf{(B)} ~18\qquad\textbf{(C)} ~24\qquad\textbf{(D)} ~32\qquad\textbf{(E)} ~44$

 \textbf{Answer: }D



\textbf{Solution 1 (Enumeration by Symmetry)}

By symmetry with respect to $3,$ note that $(x_1,x_2,x_3,x_4,x_5)$ is a valid sequence if and only if $(6-x_1,6-x_2,6-x_3,6-x_4,6-x_5)$ is a valid sequence. We enumerate the valid sequences that start with $1,2,31,$ or $32,$ as shown below:

 
                    \begin{figure}[H]
                        \centering
                        \includegraphics[width=0.5\textwidth]{Solution/AMC_10_10A_2021_20_Solution_Figure_1.png}
                    \end{figure}
                    

 There are $16$ valid sequences that start with $1,2,31,$ or $32.$ By symmetry, there are $16$ valid sequences that start with $5,4,35,$ or $34.$ So, the answer is $16+16=\boxed{\textbf{(D)} ~32}.$




\textbf{Solution 2 (Casework on the Consecutive Digits)}

Reading the terms from left to right, we have two cases for the consecutive digits, where $+$ means increase and $-$ means decrease:

 $\textbf{Case \#1: }\boldsymbol{+,-,+,-}$
$\textbf{Case \#2: }\boldsymbol{-,+,-,+}$

 For $\text{Case \#1},$ note that for the second and fourth terms, one term must be $5,$ and the other term must be either $3$ or $4.$ We have four subcases: 

 $(1) \ \underline{\hspace{3mm}}3\underline{\hspace{3mm}}5\underline{\hspace{3mm}}$
$(2) \ \underline{\hspace{3mm}}5\underline{\hspace{3mm}}3\underline{\hspace{3mm}}$
$(3) \ \underline{\hspace{3mm}}4\underline{\hspace{3mm}}5\underline{\hspace{3mm}}$
$(4) \ \underline{\hspace{3mm}}5\underline{\hspace{3mm}}4\underline{\hspace{3mm}}$

 For $(1),$ the first two blanks must be $1$ and $2$ in some order, and the last blank must be $4.$ So, we get $2$ possibilities. Similarly, $(2)$ also has $2$ possibilities.

 For $(3),$ there are no restrictions for the numbers $1, 2,$ and $3.$ So, we get $3!=6$ possibilities. Similarly, $(4)$ also has $6$ possibilities.

 Together, $\text{Case \#1}$ has $2+2+6+6=16$ possibilities. By symmetry, $\text{Case \#2}$ also has $16$ possibilities. 

 Finally, the answer is $16+16=\boxed{\textbf{(D)} ~32}.$
<u>\textbf{Remark}</u>

 This problem is somewhat similar to \href{https://artofproblemsolving.com/wiki/index.php?title=2004\_AIME\_I\_Problems/Problem\_6}{2004 AIME I Problem 6}.

 



\textbf{Solution 3 (Casework Similar to Solution 2)}

Like Solution 2, we have two cases. Due to symmetry, we just need to count one of the cases. For the purpose of this solution, we will be doing $-,+,-,+$. Instead of starting with 5, we start with 1.

 There are two ways to place it:

 \_1\_ \_ \_

 \_ \_ \_1\_

 Now we place 2, it can either be next to 1 and on the outside, or is place in where 1 would go in the other case. So now we have another two "sub case":

 \_1\_2\_(case 1)

 21\_ \_ \_(case 2)

 There are 3! ways to arrange the rest for case 1, since there is no restriction.

 For case 2, we need to consider how many ways to arrange 3,4,5 in a a>b<c fashion. It should seem pretty obvious that b has to be 3, so there will be 2! way to put 4 and 5.

 Now we find our result, times 2 for symmetry, times 2 for placement of 1 and times (3!+2!) for the two different cases for placement of 2. This give us $2*2*(3!+2!)=4*(6+2)=\boxed{\textbf{(D)} ~32}$.

 



\textbf{Solution 4 (Casework on the Position of 5)}

We only need to find the \# of rearrangements when 5 is the 4th digit and 5th digit. Find the total, and multiply by 2. Then we can get the answer by adding the case when 5 is the third digit.

 \textbf{Case $1$:} 5 is the 5th digit. \_\_ \_\_ \_\_ \_\_ 5

 Then $4$ can only be either 1st digit or the 3rd digit. 

 4 \_\_ \_\_ \_\_ 5, then the only way is that $3$ is the 3rd digit, so it can be either $231$ or $132$, give us $2$ results. 

 \_\_ \_\_ 4 \_\_ 5, then the 1st digit must be $2$ or $3$, $2$ gives us $1$ way, and $3$ gives us $2$ ways. (Can't be $1$ because the first digit would increasing). Therefore, $4$ in the middle and $5$ in the last would result in $3$ ways. 

 \textbf{Case $2$:} $5$ is the fourth digit. \_\_ \_\_ \_\_ 5 \_\_

 Then the last digit can be all of the 4 numbers $1$, $2$, $3$, and $4$. Let's say if the last digit is $4$, then the 2nd digit would be the largest for the remaining digits to prevent increasing order or decreasing order. Then the remaining two are interchangeable, give us $2!$ ways. All of the $4$ can work, so case $2$ would result in $2!+2!+2!+2!=8$ ways. 

 \textbf{Case $3$:} $5$ is in the middle. \_\_ \_\_ 5 \_\_ \_\_

 Then there are only two cases: 1. $42513$, then 4 and 3 are interchangeable, which results in $2!*2!$. Or it can be $43512$, then 4 and 2 are interchangeable, but it can not be $23514$, so there can only be 2 possible ways: $43512$, $21534$. 

 Therefore, case 3 would result in $4+2=6$ ways. 

 $8+3+2=13$, so the total ways for case 1 and case 2 with both increasing and decreasing would be $13*2=26.$

 Finally, we have $26+6=\boxed{\textbf{(D)} ~32}.$




\textbf{Solution 5 (Overcounting)}

First, we list the triples that are invalid:

 543, 542, 541, 532, 531, 521, 432, 431, 321

 By symmetry, there are the same amount of increasing triplets as there are decreasing ones. This yields 18 invalid 3 digit permutations in total.

 Suppose the triplet is ABC and the other 2 digits are X and Y. We then have 3 ways to arrange a triplet with 2 other digits.

 ABCXY, XABCY, XYABC

 X and Y can be arranged 2 ways.

 XY, YX

 This produces 18*3*2=108 permutations of invalid results. We have 5!ways to arrange 5 numbers so 120-108=12.

 Now, we must account for overcounting. For example, when 543 is counted,it only registers as one invalid permutation but in fact, it is 3 wholeinvalid permutations. We then complete this for the rest of the list:

 54321 has 543, 432, and 321

 54213 has 542 and 421

 54123 has 541 and 123

 53214 has 532 and 321

 53124 has 531 and 124

 52134 has 521 and 134

 43215 has 432 and 321

 43125 has 431 and 125

 32145 has 321 and 145

 This produces 19 values that we have overcounted but this value itself is also overcounted. We already counted 9 of the terms. This brings the final value of overcounted terms down to 10 for the decreasing triplets. By symmetry, 10 increasing triplets were overcounted.

 This gives us $120-108+20=\boxed{\textbf{(D)} ~32}.$




\textbf{Solution 6 (Using a Table)}

It is easier to consider the complement of the desired cases, so try to find the cases that DO have three integers in increasing order. First, write down the sets of three numbers that feature the numbers in increasing order. They are $123, 124, 125, 134, 135, 145, 234, 245, 345$. Each of these can be in three positions: the three are in the front with two more numbers in the back, the three are in the middle with two on each side, and two in the front and the set of three in the back. Now, count the number of combinations of $5$ numbers that each of the set of three can form that have not been previously accounted for. Also, if the set features both $3$ increasing and $3$ decreasing, then do not count it because we will separately count them. 

 
 \[\begin{tabular}[t]{|l|c|c|c|} \hline Three Increasing Numbers&Front&Middle&Back\\\hline 123&2&2&2\\\hline 124&2&2&2\\\hline 125&1&2&2\\\hline 134&2&2&2\\\hline 135&1&2&2\\\hline 145&1&2&2\\\hline 234&2&1&1\\\hline 245&1&1&1\\\hline 345&1&0&1\\\hline \end{tabular}\]

 This gives us a total of $41$ possibilities. These account for the possibilities with ONLY increasing numbers. Mirror over to the other side to get the set of combinations with either at least $1$ set of $3$ increasing numbers in a row or only at least $1$ set of $3$ decreasing numbers, but not both. We can count the ones with both separately. $14532, 12543, 13542, 54123, 53124, 52134$. Total of $6$. 

 $41\cdot 2+6=82+6=88$. This is the compliment. There are a total of $5!=120$. $120-88=32$. 

 Thus, the answer is $\boxed{\textbf{(D)} ~32}$.

 



\textbf{Solution 7 (Symmetry)}

First, note that there is a symmetry from $+,-,+,-$ to $-,+,-,+$ as follows: $a, b, c, d, e \leftrightarrow 6-a, 6-b, 6-c, 6-d, 6-e$. Now, consider the placement of 5 in the $+,-,+,-$ case. Clearly, 5 is the maximum value, so it must be placed in the 2nd position or the 4th position, but we also have symmetry $a, b, c, d, e \leftrightarrow e, d, c, b, a$. We will find the number of sequences with 5 in the second position by using casework on the position of the 4.

 \textbf{Case 1: }

 4 5 \_\_ \_\_ \_\_: 3 is the maximum value, so it must go into the 4th position, leaving us 2 ways to fill in 1 and 2.

 \textbf{Case 2: }

 \_\_ 5 \_\_ 4 \_\_: 1, 2, and 3 are all less than 4 or 5, so they can be filled in any manner: 3! = 6. 

 Overall, we have 8 ways to fill in the sequences. By our two types of symmetry, there are $8 \cdot 2 \cdot 2 = \boxed{\textbf{(D)} ~32}$.

 



\textbf{Solution 8 (PIE: No Casework)}

We use complementary counting. 

 Let $A$ be the set of all permutations $(a_1, a_2, a_3, a_4, a_5)$ of $(1,2,3,4,5)$ in which there exists at least one set of three consecutive increasing terms, and let $B$ be the set of all permutations of $(1,2,3,4,5)$ in which there exists at least one set of three consecutive decreasing terms. The complement is $|A \cup B|=|A|+|B|-|A \cap B|$. 

 First, we calculate $|A|$. Let $S_i$ be the set of permutations in which $(a_i, a_{i+1}, a_{i+2})$ is increasing for $i=1,2,3$. We have 
 \[|A|=|S_1 \cup S_2 \cup S_3|=|S_1|+|S_2|+|S_3|-|S_1 \cap S_2|-|S_2 \cap S_3|-|S_3 \cap S_1|+|S_1 \cap S_2 \cap S_3|.\]

 Since there are $\binom{5}{3}$ ways to pick $(a_1, a_2, a_3)$ (there's only one way to arrange them in increasing order) and $2!$ ways to order $a_4, a_5$, 
 \[|S_1|=\binom{5}{3} \cdot 2!=20.\] By symmetry, we find $|S_2|=|S_3|=20$. 

 Since we require $a_1<a_2<a_3<a_4$ and there are $\binom{5}{4}$ ways to pick the numbers (again, only one way to arrange them in increasing order), 
 \[|S_1 \cap S_2|=\binom{5}{4} \cdot 1!=5.\] By symmetry, $|S_2 \cap S_3|=5$.

 Notice that $S_3 \cap S_1=S_1 \cap S_2 \cap S_3$ since both translate to $a_1<a_2<a_3<a_4<a_5$, so we get 
 \[|A|=20+20+20-5-5-|S_3 \cap S_1|+|S_1\cap S_3|=60-10=50.\] By symmetry, $|B|=50$, and it remains to compute $|A\cap B|$. 

 We have two cases, either $a_1<a_2<a_3>a_4>a_5$ or $a_1>a_2>a_3<a_4<a_5$ (it is easy to verify that these are the only ones that work). WLOG consider the first. We must have $a_3=5$ since it's the largest, and for every choice of $a_1, a_2$, there is exactly one way to arrange them in order and to choose/arrange $a_3,a_4$. So, there are $\binom{4}{2}$ ways to pick $a_1, a_2$. Multiplying by 2 gives 
 \[|A\cap B|=2\cdot 6=12 \implies |A \cup B|=50+50-12=88.\]

 Since there are $5!=120$ permutations in total, the answer is $120-88=\boxed{\textbf{(D)}~32}$.

 



\textbf{Solution 9 (Quick)}

Notice that the ways to arrange the numbers is $n$, $n+2$, $n+3$, $n+2$, and $n$, for when it starts with $1$, $2$, $3$, $4$, and $5$, respectively. Now that we have found a pattern, we can just solve for $n$, so we can choose to find the ways if it starts with $1$ or $5$. Solving, we get $n$ = 5. So, the answer is $5+7+8+7+5=\boxed{\textbf{(D)}~32}$.

 


 



\textbf{Solution 10 (Brute Force)}

We write out the $5!=120$ cases, then filter out the valid ones:

 $13254,14253,14352,15243,15342,21435,21534,23154,24153,24351,25143,25341,\linebreak  31425,31524,32415,32514,34152,34251,35142,35241,41325,41523,42315,42513,\linebreak 43512,45132,45231,51324,51423,52314,52413,53412.$

 We count these out and get $\boxed{\textbf{(D)} ~32}$ permutations that work. 

 Note: We are NOT going to write our $120$ cases and factor them out. I don't know why this is the first solution to pop up.

 Note: Moved to number ten due to the sheer unnecessary nature of this solution :)

 This solution is (kind of) realistic though, as there is enough space between the solutions where you could be off by a few and still be correct.

 For anyone reading this, you shouldn't do this during the competition unless you have a lot of time and nothing else to do.~Ning2020

 



\textbf{Solution 11}

We take our $120$ cases and subtract the cases with $3$ increasing or $3$ decreasing. By symmetry, we can find the amount for one and multiply the number by $2$. 

 Originally, we find that there are $10$ choices for the numbers we will choose to be increasing, $3$ choices for where to place them, and $2$ choices for where to place the remaining numbers, for a total of $60$ ways. However, we have overcounted all of the cases with multiple triplets of increasing integers in a row. We write these out and find that we have overcounted by $10$ ways, for an actual total of $50$ ways. Be careful at this step, as for certain cases like $12345$ we actually overcounted by $2$.

 We multiply this by $2$ to find that there are also $50$ ways for the decreasing.

 However, we have now overcounted the cases in which there are increasing and decreasing triplets.

 We break these cases into $2$ different groups: the first group is where the first $3$ integers are decreasing, and the last triplet will be increasing. This forces the middle number to be $1$; we quickly write out the cases and find $6$.

 The second case is where the first $3$ are increasing and the second $3$ are decreasing. This means the middle number can be a $5$. We write out our cases for these and find another $6$ cases, for a grand total of $12$ overcounted cases. This means there are actually a total of $50+50-12=88$ distinct cases where there are $3$ increasing in a row or $3$ decreasing.

 We now subtract these $88$ cases from $120$ to get our answer of $\boxed{32}$.

 


	\item (\textit{AMC 10 10A 2021}) Let $ABCDEF$ be an equiangular hexagon. The lines $AB, CD,$ and $EF$ determine a triangle with area $192\sqrt{3}$, and the lines $BC, DE,$ and $FA$ determine a triangle with area $324\sqrt{3}$. The perimeter of hexagon $ABCDEF$ can be expressed as $m +n\sqrt{p}$, where $m, n,$ and $p$ are positive integers and $p$ is not divisible by the square of any prime. What is $m + n + p$?


 $\textbf{(A)} ~47\qquad\textbf{(B)} ~52\qquad\textbf{(C)} ~55\qquad\textbf{(D)} ~58\qquad\textbf{(E)} ~63$

 \textbf{Answer: }C



\textbf{Solution 1}

Let $P,Q,R,X,Y,$ and $Z$ be the intersections $\overleftrightarrow{AB}\cap\overleftrightarrow{CD},\overleftrightarrow{CD}\cap\overleftrightarrow{EF},\overleftrightarrow{EF}\cap\overleftrightarrow{AB},\overleftrightarrow{BC}\cap\overleftrightarrow{DE},\overleftrightarrow{DE}\cap\overleftrightarrow{FA},$ and $\overleftrightarrow{FA}\cap\overleftrightarrow{BC},$ respectively.

 The sum of the interior angles of any hexagon is $720^\circ.$ Since hexagon $ABCDEF$ is equiangular, each of its interior angles is $720^\circ\div6=120^\circ.$ By angle chasing, we conclude that the interior angles of $\triangle PBC,\triangle QDE,\triangle RFA,\triangle XCD,\triangle YEF,$ and $\triangle ZAB$ are all $60^\circ.$ Therefore, these triangles are all equilateral triangles, from which $\triangle PQR$ and $\triangle XYZ$ are both equilateral triangles.

 We are given that
 \begin{alignat*}{8} [PQR]&=\frac{\sqrt{3}}{4}\cdot PQ\textasciicircum{}2&&=192\sqrt3, \\ [XYZ]&=\frac{\sqrt{3}}{4}\cdot YZ\textasciicircum{}2&&=324\sqrt3, \end{alignat*}so we get $PQ=16\sqrt3$ and $YZ=36,$ respectively.

 By equilateral triangles and segment addition, we find the perimeter of hexagon $ABCDEF:$
 \begin{align*} AB+BC+CD+DE+EF+FA&=AZ+PC+CD+DQ+YF+FA \\ &=(YF+FA+AZ)+(PC+CD+DQ) \\ &=YZ+PQ \\ &=36+16\sqrt{3}. \end{align*}
Finally, the answer is $36+16+3=\boxed{\textbf{(C)} ~55}.$




\textbf{Solution 2}

Let the length $AB=x, BC=y.$ Then, we have
 \begin{align*} (y+2x)^2\cdot\frac{\sqrt 3}{4}&=324\sqrt3, \\ (x+2y)^2\cdot\frac{\sqrt 3}{4}&=192\sqrt3. \end{align*}
We get
 \begin{align*} y+2x&=36, \\ x+2y&=16\sqrt3. \end{align*}
We want $3x+3y,$ and it follows that 
 \[3x+3y=(y+2x)+(x+2y)=36+16\sqrt3.\]Finally, the answer is $36+16+3=\boxed{\textbf{(C)} ~55}.$

 Note: This solution is kind of a fakesolve because it assumes the hexagon is rotationally symmetric, which we can see it is not by the diagram in Solution 1.

 



\textbf{Solution 3 (Solution 1 But More Condensed)}

Since it is an equiangular hexagon, each interior angle measures $120$ degrees, and by angle chasing you can conclude that all the triangles are equilateral. Since the area formula for an equiangular triangle is $s^2 \frac{\sqrt3}{4}$, the side lengths of the triangles are $36$ and $16 \sqrt{3}$. In the above diagram the sum of the lengths of $AF$, $AB$, and $FE$ is $36$ and the sum of the lengths of $BC$, $CD$, and $DE$ is $16\sqrt{3}$. Since these are also the side lengths of the hexagon, the answer is $36+16+3=55  (C)$.

 


	\item (\textit{AMC 10 10A 2021}) Hiram's algebra notes are $50$ pages long and are printed on $25$ sheets of paper; the first sheet contains pages $1$ and $2$, the second sheet contains pages $3$ and $4$, and so on. One day he leaves his notes on the table before leaving for lunch, and his roommate decides to borrow some pages from the middle of the notes. When Hiram comes back, he discovers that his roommate has taken a consecutive set of sheets from the notes and that the average (mean) of the page numbers on all remaining sheets is exactly $19$. How many sheets were borrowed?


 $\textbf{(A)} ~10\qquad\textbf{(B)} ~13\qquad\textbf{(C)} ~15\qquad\textbf{(D)} ~17\qquad\textbf{(E)} ~20$

 \textbf{Answer: }B



\textbf{Solution 1}

Suppose the roommate took sheets $a$ through $b$, or equivalently, page numbers $2a-1$ through $2b$. Because there are $(2b-2a+2)$ numbers taken, 
 \[\frac{(2a-1+2b)(2b-2a+2)}{2}+19(50-(2b-2a+2))=\frac{50\cdot51}{2} \implies (2a+2b-39)(b-a+1)=\frac{50\cdot13}{2}=25\cdot13.\] The first possible solution that comes to mind is if $2a+2b-39=25, b-a+1=13 \implies a+b=32, b-a=12$, which indeed works, giving $b=22$ and $a=10$. The answer is $22-10+1=\boxed{\textbf{(B)} ~13}$.

 



\textbf{Solution 2}

Suppose the smallest page number borrowed is $k,$ and $n$ pages are borrowed. It follows that the largest page number borrowed is $k+n-1.$

 We have the following preconditions:

 
            \begin{enumerate} 
        \item $n$ pages are borrowed means that $\frac{n}{2}$ sheets are borrowed, from which $n$ must be even and $n<50.$

        \item $k$ must be odd, as the smallest page number borrowed is on the right side (odd-numbered).

        \item $1+2+3+\cdots+50=\frac{51(50)}{2}=1275.$

        \item The sum of the page numbers borrowed is $\frac{(2k+n-1)n}{2}.$

            \end{enumerate}
            Together, we have 
 \begin{align*} \frac{1275-\frac{(2k+n-1)n}{2}}{50-n}&=19 \\ 1275-\frac{(2k+n-1)n}{2}&=19(50-n) \\ 2550-(2k+n-1)n&=38(50-n) \\ 2550-(2k+n-1)n&=1900-38n \\ 650&=(2k+n-39)n. \end{align*}
The factors of $650$ are 
 \[1,2,5,10,13,25,26,50,65,130,325,650.\] Since $n$ is even and $n<50,$ we only have a few cases to consider:
 \[\begin{array}{c|c|c} & & \\ [-2.25ex]  \boldsymbol{n} & \boldsymbol{2k+n-39} & \boldsymbol{k} \\ [0.5ex] \hline  & & \\ [-2ex]  2 & 325 & 181 \\    10 & 65 & 47 \\  26 & 25 & 19 \end{array}\]Since $1\leq k \leq 49,$ only $k=47$ or $k=19$ is possible. If $k=47,$ then there will not be enough pages when we take $10$ pages out starting from page $47.$ Therefore, the only possibility is $k=19.$ We conclude that $n=26$ pages, or $\frac n2=\boxed{\textbf{(B)} ~13}$ sheets, are borrowed.

 



\textbf{Solution 2a (Basically the same thing as Solution 2)}

Suppose the smallest page number borrowed is $k,$ and the number of sheets is $n$. It follows that the largest page number borrowed is $k+2n-1.$

 Then the total sum of pages he took is $(2k + 2n - 1)2n/2 = (2k + 2n - 1)n$. 

 The total sum of all the pages is $\frac{50\cdot51}{2} = 1275$. 

 With the mean being 19, we can write $\frac{1275 - (2k + 2n - 1)n}{50 - 2n} = 19$. 

 Simplifying, $1275 - (2k + 2n - 1)n = 950 - 38n$ and $325 = (2k + 2n - 39)n$.

 Since $n$ (the number of sheets) is a positive integer, it must be a factor of $325 = 25\cdot13$. From the answer choices, $\boxed{\textbf{(B)} ~13}$ is the only divisor of $325$.

 -unhappyfarmer

 



\textbf{Solution 3}

Let $n$ be the number of sheets borrowed, with an average page number $k+25.5$. The remaining $25-n$ sheets have an average page number of $19$ which is less than $25.5$, the average page number of all $50$ pages, therefore $k>0$. Since the borrowed sheets start with an odd page number and end with an even page number we have $k \in \mathbb N$. We notice that $n < 25$ and $k \le (49+50)/2-25.5=24<25$.

 The weighted increase of average page number from $25.5$ to $k+25.5$ should be equal to the weighted decrease of average page number from $25.5$ to $19$, where the weights are the page number in each group (borrowed vs. remained), therefore

 
 \[2nk=2(25-n)(25.5-19)=13(25-n) \implies 13 | n \text{ or } 13 | k\]

 Since $n, k < 25$ we have either $n=13$ or $k=13$. If $n=13$ then $k=6$. If $k=13$ then $2n=25-n$ which is impossible. Therefore the answer should be $n=\boxed{\textbf{(B)} ~13}$.

 



\textbf{Solution 4}

Let $(2k-1)-2n$ be pages be borrowed, the sum of the page numbers on those pages is $(2n+2k+1)(n-k)$ while the sum of the rest pages is $1275-(2n+2k+1)(n-k)$ and we know the average of the rest is $\frac{1275-(2n+2k+1)}{50-2n+2k}$ which equals to $19$; multiply this out we got $950-38(n-k)=1275-(2n+2k+1)(n-k)$ and we got $(2n+2k-37)(n-k)=325$. As $325=25\cdot13$, we can see $n-k=13$ and that is desired $\boxed{\textbf{(B)} ~13}$.

 



\textbf{Solution 5}

Let $c$ be the number of consecutive sheets Hiram’s roommate borrows, and let $b$ be the number of sheets preceding the $c$ borrowed sheets (i.e. if the friend borrows sheets $3$, $4$, and $5$, then $c=3$ and $b=2$).

 The sum of the page numbers up till $b$ sheets is $1+2+3+\cdots + 2b=\frac{2b\cdot(2b+1)}{2} = b(2b+1)$.The last page number of the borrowed sheets would be $2(b+c)$. Therefore, the sum of the remaining page numbers of the sheets after the $c$ borrowed sheets would be $2(b+c)+1 + 2(b+c)+2+\cdots+50$.

 The total number of page numbers after the borrow would be $50-2c$.

 Thus the average of the page numbers after the borrow would be:
 \[\frac{b(2b+1)+ 2(b+c)+1 + 2(b+c)+2+\cdots+50}{50-2c} =19.\]By the arithmetic series formula, this turns out to be:
 \[\frac{b(2b+1)+ \frac{(2(b+c)+1+50)\cdot(50-2c-2b)}{2}}{50-2c} =19\]because in the changed sum, there are $50$ numbers minus $2c$ borrowed numbers and $2b$ numbers from the first $b$ sheets.

 This simplifies to 
 \[\frac{b(2b+1)+ (2(b+c)+51)\cdot(25-c-b)}{50-2c} =19 \implies 19(50-2c)= b(2b+1)+ (2b+2c+51)\cdot(25-c-b).\]Noticing that some terms will cancel, we expand, leading to:
 \[950-38c=2b^2+b+50b-2bc-2b^2+50c-2c^2-2bc+1275-51c-51b\]
 \[\implies 950-38c=1275-4bc-c-2c^2 \implies 2c^2+4b-37c=325.\]Factoring, we get 
 \[c(2c+4b-37)=325.\]The prime factorization of 325 is $5^2\cdot13$. Recall that $0\leq c \leq 25$, so $c$ could be $1$, $5$, $13$, or $25$.

 We can rule out $c=25$ since Hiram would have no paper left over, so the average of the page numbers he has would be $0$. We can now plug in the other answers for $c$ and we if we get a valid answer for $b$.

 
            \begin{itemize}
        \item $c=1 \implies 4b-35=325 \implies 4b=360 \implies b=90$. Since Hiram only has $25$ sheets, this is clearly wrong and we can rule out $c=1$.
            \end{itemize}
            
            \begin{itemize}
        \item $c=5 \implies 10+4b-37=65 \implies 4b = 92 \implies b=23$. Since $b+c$ is $28$ but we have only $25$ sheets, this is also implausible so we can rule out $c=5$.
            \end{itemize}
            Finally, just to make sure, we test $c=13 \implies 26+4b-37=25 \implies 4b=36 \implies b=9$. 

 $b$ is an integer and $b+c=22$, so everything checks out. The number of consecutive sheets borrowed by Hiram’s friend is $c=\boxed{\textbf{(B)} ~13}$.

 



\textbf{Solution 6}

The sum of all the page numbers is
 \[1+2+3+\cdots+50 = 1275.\]If we add the page numbers on each sheet, we get this sequence:
 \[3, 7, 11, \ldots, 99.\]So we can write the sum of the numbers on the first sheet that the roommate borrowed as $4n+3$ for some nonnegative integer, $n$. If the roommate borrowed $k$ sheets, he borrowed sheets
 \[4n+3, 4n+7, \ldots, 4n+4k-1.\]The sum of the numbers in this sequence is
 \[\frac{k(8n+4k+2)}{2} = 4nk+2k^2+k.\]Since there are $2$ pages per sheet, there are $50-2k$ pages remaining, so the average page number of the remaining sheets is 
 \[\frac{1275 - (4nk+2k^2+k)}{50-2k}.\]Therefore,
 \[\frac{1275 - (4nk+2k^2+k)}{50-2k} = 19,\]which simplifies to
 \[2k^2-37k-325 = -4nk.\]Factoring the left-hand side,
 \[(2k+13)(k-25)=-4nk.\] Since the right-hand side of this equation is divisible by $k$, the left-hand side must also be divisible by $k$. 

 In order for $(2k+13)(k-25)$ to be divisible by $k$, either $2k+13$ or $k-25$ must be divisible by $k$. 
 \[\frac{2k+13}{k}=2+\frac{13}{k},\] so $2k+13$ is divisible by $k$ only if $k$ is a factor of $13$. 
 \[\frac{k-25}{k}=1-\frac{25}{k},\] so $k-25$ is divisible by $k$ only if $k$ is a factor of $25$. 

 None of the answer choices are factors of $25$, but answer choice B is a factor of $13$. Hence, the answer is $\boxed{\textbf{(B)} ~13}$.

 



\textbf{Solution 7}

We see the sum of the first $50$ page numbers is $1275$. We can then write an equation to solve for $k$, the number of pages removed. Let the average of the page numbers removed be $a$. We write $\frac{1275-ka}{50-k}=19$. 

 We simplify this equation to get $325=k(a-19)$.

 We now test multiples of $325$.

 We know that $a$ can be at max $49$, so $k=1$ and $k=5$ won't work.

 We test $k=13$, and this value works. $k=25$ also works, but this is not an answer choice, so we see that the answer must be $\boxed{13}$.

 


	\item (\textit{AMC 10 10A 2021}) Frieda the frog begins a sequence of hops on a $3 \times 3$ grid of squares, moving one square on each hop and choosing at random the direction of each hop-up, down, left, or right. She does not hop diagonally. When the direction of a hop would take Frieda off the grid, she "wraps around" and jumps to the opposite edge. For example if Frieda begins in the center square and makes two hops "up", the first hop would place her in the top row middle square, and the second hop would cause Frieda to jump to the opposite edge, landing in the bottom row middle square. Suppose Frieda starts from the center square, makes at most four hops at random, and stops hopping if she lands on a corner square. What is the probability that she reaches a corner square on one of the four hops?


 $\textbf{(A)} ~\frac{9}{16}\qquad\textbf{(B)} ~\frac{5}{8}\qquad\textbf{(C)} ~\frac{3}{4}\qquad\textbf{(D)} ~\frac{25}{32}\qquad\textbf{(E)} ~\frac{13}{16}$

 \textbf{Answer: }D



\textbf{Solution 1 (Complementary Counting)}

We will use complementary counting.First, the frog can go left with probability $\frac14$. We observe symmetry, so our final answer will be multiplied by 4 for the 4 directions, and since $4\cdot \frac14=1$, we will ignore the leading probability.

 From the left, she either goes left to another edge ($\frac14$) or back to the center ($\frac14$). Time for some casework.

 $\textbf{Case 1:}$ She goes back to the center.

 Now, she can go in any 4 directions, and then has 2 options from that edge. This gives $\frac12$. --End case 1

 $\textbf{Case 2:}$ She goes to another edge (rightmost).

 Subcase 1: She goes back to the left edge. She now has 2 places to go, giving $\frac12$

 Subcase 2: She goes to the center. Now any move works.

 $\frac14 \cdot \frac12 + \frac14 \cdot 1=\frac18 + \frac 14=\frac38$ for this case. --End case 2

 She goes back to the center in Case 1 with probability $\frac14$, and to the right edge with probability $\frac14$

 So, our answer is $\frac14 \cdot \frac12 + \frac14 \cdot \frac38=\frac14 (\frac12+\frac38)=\frac14 \cdot \frac78 = \frac7{32}$

 But, don't forget complementary counting. So, we get $1-\frac7{32}=\frac{25}{32} \implies \boxed{D}$. ~ firebolt360

 



\textbf{Solution 2 (Direct Counting and Probability States)}

We can draw a state diagram with three states: center, edge, and corner. Denote center by M, edge by E, and corner by C. There are a few ways Frieda can reach a corner in four or less moves: EC, EEC, EEEC, EMEC. Then, calculating the probabilities of each of these cases happening, we have $1\cdot\frac{1}{2}+1\cdot\frac{1}{4}\cdot\frac{1}{2}+1\cdot\frac{1}{4}\cdot\frac{1}{4}\cdot\frac{1}{2}+1\cdot\frac{1}{4}\cdot1\cdot\frac{1}{2}=\frac{25}{32}$, so the answer is $\boxed{D}$. ~IceWolf10

 



\textbf{Solution 3 (Finds Numerator and Denominator Separately)}

Suppose Frieda makes four independent hops without stopping so that each outcome is equally likely.

 <u>\textbf{Denominator}</u>

 There are $4^4=256$ ways for Frieda to make four hops without restrictions.

 <u>\textbf{Numerator}</u>

 We perform casework on which hop Frieda reaches a corner for the first time:

 
            \begin{enumerate} 
        \item The second hop (The third and the fourth hops have no restrictions.)The four hops have $4, 2, 4, 4$ options, respectively. So, this case has $4\cdot2\cdot4\cdot4=128$ ways. 

  

 
        \item The third hop (The fourth hop has no restrictions.)No matter which direction the first hop takes, the second hop must "wrap around". 

 The four hops have $4, 1, 2, 4$ options, respectively. So, this case has $4\cdot1\cdot2\cdot4=32$ ways. 

  

 
        \item The fourth hop (There are two subcases based on the second hop.)
<ol style="margin-left: 1.5em;" type="A"> 
        \item The second hop "wraps around". It follows that the third hop must also "wrap around".The four hops have $4, 1, 1, 2$ options, respectively. So, this subcase has $4\cdot1\cdot1\cdot2=8$ ways. 

  

 
        \item The second hop returns to the center.The four hops have $4, 1, 4, 2$ options, respectively. So, this subcase has $4\cdot1\cdot4\cdot2=32$ ways.     

 </ol>Together, this case has $8+32=40$ ways.

 
            \end{enumerate}
            The numerator is $128+32+40=200.$
<u>\textbf{Probability}</u>

 The requested probability is $\frac{200}{256}=\boxed{\textbf{(D)} ~\frac{25}{32}}.$
<u>\textbf{Remark}</u>

 This problem is quite similar to \href{https://artofproblemsolving.com/wiki/index.php?title=1995\_AIME\_Problems/Problem\_3}{1995 AIME Problem 3}.

 



\textbf{Solution 4}

Let $C_n$ be the probability that Frieda is on the central square after n moves, $E_n$ be the probability that Frieda is on one of the four squares on the middle of the edges after n moves, and $V_n$ (V for vertex) be the probability that Frieda is on a corner after n moves. The only way to reach the center is by moving in $1$ specific direction out of $4$ total directions from the middle of an edge, so $C_{n+1}=\frac{E_n}{4}$. The ways to reach the middle of an edge are by moving in any direction from the center or by moving in $1$ specific direction from the middle of an edge, so $E_{n+1}=C_n+\frac{E_n}{4}$. The ways to reach a corner are by simply staying there after reaching there in a previous move or by moving in $2$ specific directions from the middle of an edge, so $V_{n+1}=V_n+\frac{E_n}{2}$. Since Frieda always start from the center, $C_0=1$, $E_0=0$, and $V_0=0$. We use the previous formulas to work out $V_4$ and find it to be $\boxed{\textbf{(D)} ~\frac{25}{32}}$.

 -SmileKat32

 



\textbf{Solution 5}

Imagine an infinite grid of $2$ by $2$ squares such that there is a $2$ by $2$ square centered at $(3x, 3y)$ for all ordered pairs of integers $(x, y).$

                    \begin{figure}[H]
                        \centering
                        \includegraphics[width=0.5\textwidth]{Solution/AMC_10_10A_2021_23_Solution_Figure_1.png}
                    \end{figure}
                    

 It is easy to see that the problem is equivalent to Frieda moving left, right, up, or down on this infinite grid starting at $(0, 0)$ (minus the teleportations). Since counting the complement set is easier, we'll count the number of $4$-step paths such that Frieda never reaches a corner point. 

 In other words, since the reachable corner points are $(\pm 1, \pm 1), (\pm 1, \pm 2), (\pm 2, \pm 1),$ and $(\pm 2, \pm 2),$ Frieda can only travel along the collection of points included in $S$, where $S$ is all points on $x=0$ and $y=0$ such that $|y|<4$ and $|x|<4$, respectively, plus all points on the big square with side length $6$ centered at $(0, 0).$ We then can proceed with casework:

 Case $1$: Frieda never reaches $(0, \pm 3)$ nor $(\pm 3, 0).$

 When Frieda only moves horizontally or vertically for her four moves, she can do so in $2^4 - 4 = 12$ ways for each case . Thus, $12 \cdot 2$ total paths for the subcase of staying in one direction. (For instance, all length $4$ combinations of $F$ and $B$ except $FFFF$, $BBBB$, $FFFB$, and $BBBF$ for the horizontal direction.) 

 There is another subcase where she changes directions during her path. There are four symmetric cases for this subcase depending on which of the four quadrants Frieda hugs. For the first quadrant, the possible paths are $FBUD$, $FBUU$, $UDFB$, and $UDFF.$ Thus, a total of $4 \cdot 4 = 16$ ways for this subcase. 

 Total for Case $1$: $24 + 16 = 40$

 Case $2$: Frieda reaches $(0, \pm 3)$ or $(\pm 3, 0)$.

 Once Frieda reaches one of the points listed above (by using three moves), she has four choices for her last move. Thus, a total of $4 \cdot 4 = 16$ paths for this case.

 Our total number of paths never reaching coroners is thus $16+40=56,$ making for an answer of 
 \[\frac{4^4-56}{4^4} = \boxed{\textbf{(D)} ~\frac{25}{32}}.\]

 -fidgetboss\_4000

 



\textbf{Solution 6 (Casework)}

We take cases on the number of hops needed to reach a corner. For simplicity, denote $E$ as a move that takes Frieda to an edge, $W$ as wrap-around move and $C$ as a corner move. Also, denote $O$ as a move that takes us to the center.

 \textbf{ 2 Hops }

 Then, Frieda will have to $(E, C)$ as her set of moves. There are $4$ ways to move to an edge, and $2$ corners to move to, for a total of $4 \cdot 2 = 8$ cases here. Then, there are $4$ choices for each move, for a probability of $\frac{8}{4 \cdot 4} = \frac{1}{2}$. 

 \textbf{ 3 Hops }

 In this case, Frieda must wrap-around. There's only one possible combination, just $(E, W, C)$. There are $4$ ways to move to an edge, $1$ way to wrap-around (you must continue in the same direction) and $2$ corners, for a total of $4 \cdot 1 \cdot 2 = 8$ cases here. Then, there are $4$ choices for each move, for a probability of $\frac{8}{4 \cdot 4 \cdot 4} = \frac{1}{8}$. 

 \textbf{ 4 Hops } 

 Lastly, there are two cases we must consider here. The first case is $(E, O, E, C)$, and the second is $(E, W, W, C)$. For the first case, there are $4$ ways to move to an edge, $1$ way to return to the center, $4$ ways to move to an edge once again, and $2$ ways to move to a corner. Hence, there is a total of $4 \cdot 1 \cdot 4 \cdot 2 = 32$ cases here. Then, for the second case, there are $4$ ways to move to a corner, $1$ way to wrap-around, $1$ way to wrap-around the way Frieda just came from, and $2$ ways to move to a corner. This implies there are $4 \cdot 1 \cdot 1 \cdot 2 = 8$ cases here. Then, there is a total of $8+32 = 40$ cases, out of a total of $4^4 = 256$ cases, for a probability of $\frac{40}{256} = \frac{5}{32}$. 

 Then, the total probability that Frieda ends up on a corner is $\frac{1}{2} + \frac{1}{8} + \frac{5}{32} = \frac{25}{32}$, corresponding to choice $\boxed{\textbf{(D)}  ~\frac{25}{32}}$. 

 



\textbf{Solution 7}

I denote 3x3 grid by

 - HOME square (x1)

 - CORN squares (x4)

 - SIDE squares (x4)

 Transitions:

 - HOME always move to SIDE

 - CORN is DONE

 - SIDE move to HOME with $p=1/4,$ move to SIDE with $p=1/4,$ and move to CORN with $p=1/2.$

 After one move, will be on $\text{SIDE}$ square

 After two moves, will be $1/2 + 1/4 (\text{HOME}) + 1/4 (\text{SIDE})$

 After three moves, will be $1/2 + 1/4 (\text{SIDE}) + 1/4 (1/2 + 1/4(\text{SIDE}) + 1/4(\text{HOME}))$

 After four moves, probability on CORN will be $1/2 + 1/4 (1/2) + 1/4( 1/2 + 1/4(1/2)) = 1/2 + 1/8 + 5/32 = 25 / 32.$




\textbf{Solution 8 (Markov Chain)}


            \begin{enumerate}[label=(\roman*)]
                \item <i>Main article: \href{https://artofproblemsolving.com/wiki/index.php?title=Markov\_Chains}{Markov Chains}</i>
            \end{enumerate}
            The grid is symmetrical. Denote the center cell as state Center, the $4$ middle side cells as state Side, the $4$ corner cells as state Corner. We can draw the following State Transition Diagram with \href{https://en.wikipedia.org/wiki/Markov\_chain}{Markov Chain}. The numbers on the transition arc are the transition probability. 

 
                    \begin{figure}[H]
                        \centering
                        
\includegraphics[width=0.5\textwidth]{Solution/AMC_10_10A_2021_23_Solution_Figure_2.png}
                    \end{figure}
                    We calculate the probability of different states in each round in the following table by iterating a few more times:
 \[\begin{array}{|c|c|c|c|} \hline \textbf{Round} & \textbf{Center} & \textbf{Side} & \textbf{Corner}\\ \hline &&&\\ 1 & 0 & 1 & 0\\ &&&\\ \hline &&&\\ 2 & \frac{1}{4} & \frac{1}{4} & \frac{1}{2}\\ &&&\\ \hline &&&\\ 3 & \frac{1}{4} \cdot \frac{1}{4} = \frac{1}{16} & \frac{1}{4} + \frac{1}{4} \cdot \frac{1}{4} = \frac{5}{16} & \frac{1}{4} \cdot \frac{1}{2} + \frac{1}{2} = \frac{5}{8}\\ &&&\\ \hline &&&\\ 4 & \frac{5}{16} \cdot \frac{1}{4} = \frac{5}{64} & \frac{1}{16} + \frac{5}{16} \cdot \frac{1}{4} = \frac{9}{64} & \frac{5}{16} \cdot \frac{1}{2} + \frac{5}{8} = \frac{25}{32}\\ &&&\\ \hline \end{array}\]

 Therefore, the answer is $\boxed{\textbf{(D)}  ~\frac{25}{32}}$.

 



\textbf{Solution 9}

Let us go through the frog’s jumps step by step to find the probability that it will land in a corner. Denote the frog’s position as $F$:
                    \begin{figure}[H]
                        \centering
                        \includegraphics[width=0.5\textwidth]{Solution/AMC_10_10A_2021_23_Solution_Figure_3.png}
                    \end{figure}
                    First notice that since the frog only moves up, down, left or right, it’s original hop will not matter by symmetry, because all first hops will bring the frog at the edge next to two corners. WLOG we let the frog hop up:
                    \begin{figure}[H]
                        \centering
                        \includegraphics[width=0.5\textwidth]{Solution/AMC_10_10A_2021_23_Solution_Figure_4.png}
                    \end{figure}
                    Notice that in its second hop, there are two ways for it to move to a corner (left and right) out of four, and thus the probability of Frieda reaching a corner on the second hop is $\frac{1}{2}$.

 There are two other cases for Frieda’s second hop: it could either wrap around to another edge square or move back to the middle. We analyze each of these cases individually, noting that each has probability $\frac{1}{4}$:
                    \begin{figure}[H]
                        \centering
                        \includegraphics[width=0.5\textwidth]{Solution/AMC_10_10A_2021_23_Solution_Figure_5.png}
                    \end{figure}
                    In this case, there are again $2$ choices to go to a corner on the third hop, go back to the previous square, or go to the middle. Here there is a $\frac{1}{2}$ probability that Frieda will reach a corner on the third hop. If Frieda hops back with a $\frac{1}{4}$ probability, she will again have a $\frac{1}{2}$ chance of reaching a corner on her fourth hop, but after four hops she stops hopping, so there are no more chances for her to reach a corner if she hops back. If she hops to the middle, then her fourth hop cannot reach a corner and that case yields $0$ probability of reaching a corner.

 To sum up this case, the total probability of her reaching a corner is $\frac{1}{4}\left(\frac{1}{2}+\frac{1}{4}\cdot \frac{1}{2}\right)$.

 Finally suppose Frieda goes back to her initial position after $2$ steps. Then her third step again doesn’t matter because whatever she does, she will be next to $2$ corners, so the probability of her reaching a corner on her fourth hop in this case is$\frac{1}{4}\cdot\frac{1}{2}$.

 Summing up all the cases, we get 
 \[\frac{1}{2}+\frac{1}{4}\left(\frac{1}{2}+\frac{1}{4}\cdot\frac{1}{2}\right)+\frac{1}{4}\cdot\frac{1}{2} = \frac{1}{2}+\frac{1}{4}\left(\frac{1}{2}+\frac{1}{8}\right)+\frac{1}{8} = \frac{5}{4}\cdot\left(\frac{1}{2}+\frac{1}{8}\right) = \frac{5}{4} \cdot \frac{5}{8} = \boxed{\textbf{(D)} ~\frac{25}{32}}.\]~KingRavi

 



\textbf{Solution 10 (Similar to the States Solution)}

Denote the center square as $C$, the edge squares as $E$, and the corner squares as $Co$. On the first move, the probability of Frieda on the $C$ is $0$ because she's moving away from the center square. The probability she's on $E$ is $1$ because no matter which direction she goes, she's always going to end up on an edge square. Probability she's on $Co$ equals $1 - (0 + 1) = 0$ by complementary counting. Notice that the probability Frieda will be on the $C$ can be expressed as: $P(C) = \frac{1}{4} \cdot P(E from last move)$ because one-fourth of the moves Frieda takes from an edge square sends her back to $C$. Also, she can't reach $C$ from a $Co$ square. Now $P(E) = P(C from last move) + \frac{1}{4} \cdot P(E from last move)$ simply from the same logic. We can find $P(Co)$ using complementary counting. On the second move, $P(C) = \frac{1}{4} \cdot 1 = \frac{1}{4}$ and $P(E) = 0 + \frac{1}{4} \cdot 1 = \frac{1}{4}$. Now $P(Co) = 1 - \frac{1}{4} - \frac{1}{4} = \frac{1}{2}$. On the third move, $P(C) = \frac{1}{4} \cdot \frac{1}{4} = \frac{1}{16}$ and $P(E) = \frac{1}{4} + \frac{1}{4} \cdot \frac{1}{4} = \frac{5}{16}$. $P(Co) = 1 - \frac{1}{16} - \frac{5}{16} = \frac{5}{8}$. Now on the fourth and last move, $P(C) = \frac{1}{4} \cdot \frac{5}{16} = \frac{5}{64}$ and $P(E) = \frac{1}{16} + \frac{1}{4} \cdot \frac{5}{16} = \frac{9}{64}$. Finally, $P(Co) = 1 - \frac{5}{64} - \frac{9}{64} = 1 - \frac{7}{32} = \frac{25}{32}$.

 


	\item (\textit{AMC 10 10A 2021}) The interior of a quadrilateral is bounded by the graphs of $(x+ay)^2 = 4a^2$ and $(ax-y)^2 = a^2$, where $a$ is a positive real number. What is the area of this region in terms of $a$, valid for all $a > 0$?


 $\textbf{(A)} ~\frac{8a^2}{(a+1)^2}\qquad\textbf{(B)} ~\frac{4a}{a+1}\qquad\textbf{(C)} ~\frac{8a}{a+1}\qquad\textbf{(D)} ~\frac{8a^2}{a^2+1}\qquad\textbf{(E)} ~\frac{8a}{a^2+1}$

 \textbf{Answer: }D



\textbf{Solution 1 (Generalized Value of a)}

The cases for $(x+ay)^2 = 4a^2$ are $x+ay = \pm2a,$ or two parallel lines. We rearrange each case and construct the table below:
 \[\begin{array}{c||c|c|c|c} & & & & \\ [-2.5ex] \textbf{Case} & \textbf{Line's Equation} & \boldsymbol{x}\textbf{-Intercept} & \boldsymbol{y}\textbf{-Intercept} & \textbf{Slope} \\ [0.5ex] \hline & & & & \\ [-1.5ex] 1 & x+ay-2a=0 & 2a & 2 & -\frac1a  \\ [2ex]  2 & x+ay+2a=0 & -2a & -2 & -\frac1a \\ [0.75ex] \end{array}\]The cases for $(ax-y)^2 = a^2$ are $ax-y=\pm a,$ or two parallel lines. We rearrange each case and construct the table below:
 \[\begin{array}{c||c|c|c|c} & & & & \\ [-2.5ex] \textbf{Case} & \textbf{Line's Equation} & \boldsymbol{x}\textbf{-Intercept} & \boldsymbol{y}\textbf{-Intercept} & \textbf{Slope} \\ [0.5ex] \hline & & & & \\ [-1.5ex] 1* & ax-y-a=0 & 1 & -a & a  \\ [2ex]  2* & ax-y+a=0 & -1 & a & a \\ [0.75ex] \end{array}\]Since the slopes of intersecting lines $(1)\cap(1*), (1)\cap(2*), (2)\cap(1*),$ and $(2)\cap(2*)$ are negative reciprocals, we get four right angles, from which the quadrilateral is a rectangle.

 Two solutions follow from here:

 



\textbf{Solution 1.1 (Distance Between Parallel Lines)}

Recall that for constants $A,B,C_1$ and $C_2,$ the distance $d$ between parallel lines$\begin{cases} Ax+By+C_1=0 \\ Ax+By+C_2=0 \end{cases}$ is 
 \[d=\frac{\left|C_2-C_1\right|}{\sqrt{A^2+B^2}}.\]From this formula:

 
            \begin{itemize}
        \item The distance between lines $(1)$ and $(2)$ is $\frac{4a}{\sqrt{1+a^2}},$ the length of this rectangle.
            \end{itemize}
            
            \begin{itemize}
        \item The distance between lines $(1*)$ and $(2*)$ is $\frac{2a}{\sqrt{a^2+1}},$ the width of this rectangle.
            \end{itemize}
            The area we seek is 
 \[\frac{4a}{\sqrt{1+a^2}}\cdot\frac{2a}{\sqrt{a^2+1}}=\boxed{\textbf{(D)} ~\frac{8a^2}{a^2+1}}.\]~MRENTHUSIASM

 



\textbf{Solution 1.2 (Distance Between Points)}

The solutions to systems of equations $(1)\cap(1*), (1)\cap(2*), (2)\cap(2*), (2)\cap(1*)$ are 
 \[(x,y)=\left(\frac{a(a+2)}{a^2+1},\frac{a(2a-1)}{a^2+1}\right), \left(-\frac{a(a-2)}{a^2+1},\frac{a(2a+1)}{a^2+1}\right), \left(-\frac{a(a+2)}{a^2+1},-\frac{a(2a-1)}{a^2+1}\right), \left(\frac{a(a-2)}{a^2+1},-\frac{a(2a+1)}{a^2+1}\right),\] respectively, which are the consecutive vertices of this rectangle.

 By the Distance Formula, the length and width of this rectangle are $\frac{4a\sqrt{a^2+1}}{a^2+1}$ and $\frac{2a\sqrt{a^2+1}}{a^2+1},$ respectively.

 The area we seek is 
 \[\frac{4a\sqrt{a^2+1}}{a^2+1}\cdot\frac{2a\sqrt{a^2+1}}{a^2+1}=\boxed{\textbf{(D)} ~\frac{8a^2}{a^2+1}}.\]~MRENTHUSIASM

 



\textbf{Solution 2 (Specified Value of a)}

In this solution, we will refer to equations $(1),(2),(1*),$ and $(2*)$ from Solution 1.

 Substituting $a=2$ into the answer choices gives 

 $\textbf{(A)} ~\frac{32}{9}\qquad\textbf{(B)} ~\frac{8}{3}\qquad\textbf{(C)} ~\frac{16}{3}\qquad\textbf{(D)} ~\frac{32}{5}\qquad\textbf{(E)} ~\frac{16}{5}$

 At $a=2,$ the solutions to systems of equations $(1)\cap(1*), (1)\cap(2*), (2)\cap(2*), (2)\cap(1*)$ are 
 \[(x,y)=\left(\frac 85, \frac 65\right), (0,2), \left(-\frac 85, -\frac 65\right), (0,-2),\] respectively, which are the consecutive vertices of the quadrilateral.

 Two solutions follow from here:

 



\textbf{Solution 2.1 (Area of a Rectangle)}

From the tables in Solution 1, we conclude that the quadrilateral is a rectangle.

 By the Distance Formula, the length and width of this rectangle are $\frac{8\sqrt5}{5}$ and $\frac{4\sqrt5}{5},$ respectively.

 The area we seek is 
 \[\frac{8\sqrt5}{5}\cdot\frac{4\sqrt5}{5}=\frac{32}{5},\] from which the answer is $\boxed{\textbf{(D)} ~\frac{8a^2}{a^2+1}}.$




\textbf{Solution 2.2 (Area of a General Quadrilateral)}

Even if we do not recognize that the quadrilateral is a rectangle, we can apply the Shoelace Theorem to its <i>\textbf{consecutive}</i> vertices
 \begin{align*} (x_1,y_1) &= \left(\frac 85, \frac 65\right), \\  (x_2,y_2) &= (0,2), \\  (x_3,y_3) &= \left(-\frac 85, -\frac 65\right), \\  (x_4,y_4) &= (0,-2). \end{align*}
The area we seek is 
 \[\frac{1}{2} \left|(x_1y_2 + x_2y_3 + x_3y_4 + x_4y_1) - (y_1x_2 + y_2x_3 + y_3x_4 + y_4x_1)\right| = \frac{32}{5}.\] from which the answer is $\boxed{\textbf{(D)} ~\frac{8a^2}{a^2+1}}.$




\textbf{Solution 3 (Slopes and Intercepts)}


                    \begin{figure}[H]
                        \centering
                        \includegraphics[width=0.5\textwidth]{Solution/AMC_10_10A_2021_24_Solution_Figure_1.png}
                    \end{figure}
                    The quadrilateral is enclosed by four lines. Similar to Solution 1, we will use the equations from the four cases:

 
            \begin{enumerate} 
        \item $x+ay=2a.$ This is a line with $x$-intercept $2a,$ $y$-intercept $2,$ and slope $-\frac 1a.$

        \item $x+ay=-2a.$ This is a line with $x$-intercept $-2a,$ $y$-intercept $-2,$ and slope $-\frac 1a.$

        \item $ax-y=a.$ This is a line with $x$-intercept $1,$ $y$-intercept $-a,$ and slope $a.$

        \item $ax-y=-a.$ This is a line with $x$-intercept $-1,$ $y$-intercept $a,$ and slope $a.$

            \end{enumerate}
            It follows that $DF = 4$ and $DE = \sqrt{4^2 - s^2}$.

 Because the slope of line $y = -\frac{x}{a} + 2$ is $-\frac{1}{a}$, $\frac{1}{a} = \frac{DE}{EF} = \frac{\sqrt{16-s^2}}{s}$, $s^2(a^2+1) = 16a^2$, $s = \frac{4a}{\sqrt{a^2+1}}$.

 It follows that $AC = 2a$ and $BC = \sqrt{(2a)^2 - w^2}$.

 Because the slope of line $y = ax - a$ is $a$, $a = \frac{BC}{AB} = \frac{\sqrt{4a^2-w^2}}{w}$, $w^2(a^2+1)=4a^2$, $w=\frac{2a}{\sqrt{a^2+1}}$.

 Therefore, the answer is 
 \[\text{Area} = s \cdot w=\frac{4a}{\sqrt{a^2+1}} \cdot \frac{2a}{\sqrt{a^2+1}} = \boxed{\textbf{(D)} ~\frac{8a^2}{a^2+1}}.\]




\textbf{Solution 4 (Trigonometry)}

Similar to Solution 1, we will use the equations from the four cases:

 
            \begin{enumerate} 
        \item $x+ay=2a.$ This is a line with $x$-intercept $2a,$ $y$-intercept $2,$ and slope $-\frac 1a.$

        \item $x+ay=-2a.$ This is a line with $x$-intercept $-2a,$ $y$-intercept $-2,$ and slope $-\frac 1a.$

        \item $ax-y=a.$ This is a line with $x$-intercept $1,$ $y$-intercept $-a,$ and slope $a.$

        \item $ax-y=-a.$ This is a line with $x$-intercept $-1,$ $y$-intercept $a,$ and slope $a.$

            \end{enumerate}
            Let $\tan A=a.$ The area of the rectangle created by the four equations can be written as 
 \begin{align*} 2a\cdot \cos A\cdot4\sin A &= 8a\cos A \cdot \sin A \\ &= 8a\cdot~\frac{1}{\sqrt{a^2+1}}\cdot~\frac{a}{\sqrt{a^2+1}} \\ &= \boxed{\textbf{(D)} ~\frac{8a^2}{a^2+1}}. \end{align*}
~fnothing4994 (Solution)

 



\textbf{Solution 5 (Observations)}

The conditions $(x+ay)^2 = 4a^2$ and $(ax-y)^2 = a^2$ give $|x+ay| = |2a|$ and $|ax-y| = |a|$ or $x+ay = \pm 2a$ and $ax-y = \pm a$. The slopes here are perpendicular, so the quadrilateral is a rectangle.Plug in $a=1$ and graph it. We quickly see that the area is $2\sqrt{2} \cdot \sqrt{2} = 4$, so the answer can't be $\textbf{(A)}$ or $\textbf{(B)}$ by testing the values they give (test it!). Now plug in $a=2$. We see using a ruler that the sides of the rectangle are about $\frac74$ and $\frac72$. So the area is about $\frac{49}8 = 6.125$. Testing $\textbf{(C)}$, we get $\frac{16}3$ which is clearly less than $6$, so it is out. Testing $\textbf{(D)}$, we get $\frac{32}5$ which is near our answer, so we leave it. Testing $\textbf{(E)}$, we get $\frac{16}5$, way less than $6$, so it is out. So, the only plausible answer is $\boxed{\textbf{(D)} ~\frac{8a^2}{a^2+1}}$.

 



\textbf{Solution 6 (Observations)}

Trying $a = 1$ narrows down the choices to options $\textbf{(C)}$, $\textbf{(D)}$ and $\textbf{(E)}$. Trying $a = 2$ and $a = 3$ eliminates $\textbf{(C)}$ and $\textbf{(E)}$, to obtain $\boxed{\textbf{(D)} ~\frac{8a^2}{a^2+1}}$ as our answer. Refer to Solution 2 for a detailed explanation.

 



\textbf{Solution 7 (Observations: Cheap)}

Note that $a=2$ yields different values for all answer choices. If we put in $a=2,$ we find that the area of the quadrilateral is $\frac{32}{5}.$ This means that the answer must be $\boxed{\textbf{(D)} ~\frac{8a^2}{a^2+1}}.$ Refer to Solution 2 for a detailed explanation.

 


	\item (\textit{AMC 10 10A 2021}) How many ways are there to place $3$ indistinguishable red chips, $3$ indistinguishable blue chips, and $3$ indistinguishable green chips in the squares of a $3 \times 3$ grid so that no two chips of the same color are directly adjacent to each other, either vertically or horizontally?


 $\textbf{(A)} ~12\qquad\textbf{(B)} ~18\qquad\textbf{(C)} ~24\qquad\textbf{(D)} ~30\qquad\textbf{(E)} ~36$

 \textbf{Answer: }E



\textbf{Solution 1 (Casework on the Center's Color Chip's Configurations)}

Call the different colors A,B,C. There are $3!=6$ ways to rearrange these colors to these three letters, so $6$ must be multiplied after the letters are permuted in the grid. WLOG assume that A is in the center. 
 \[\begin{tabular}{ c c c }  ? & ? & ? \\   ? & A & ? \\    ? & ? & ?     \end{tabular}\]In this configuration, there are two cases, either all the A's lie on the same diagonal:
 \[\begin{tabular}{ c c c }  ? & ? & A \\   ? & A & ? \\    A & ? & ?     \end{tabular}\]or all the other two A's are on adjacent corners:
 \[\begin{tabular}{ c c c }  A & ? & A \\   ? & A & ? \\    ? & ? & ?     \end{tabular}\]In the first case there are two ways to order them since there are two diagonals, and in the second case there are four ways to order them since there are four pairs of adjacent corners.

 In each case there is only one way to put the three B's and the three C's as shown in the diagrams. 
 \[\begin{tabular}{ c c c }   C & B & A \\    B & A & C \\     A & C & B     \end{tabular}\]
 \[\begin{tabular}{ c c c }  A & B & A \\   C & A & C \\    B & C & B     \end{tabular}\]This means that there are $4+2=6$ ways to arrange A,B, and C in the grid, and there are $6$ ways to rearrange the colors. Therefore, there are $6\cdot6=36$ ways in total, which is $\boxed{\textbf{(E)} ~36}$.

 -happykeeper

 



\textbf{Solution 2 (Casework on the Top-Center and Center-Left Chips)}

Without the loss of generality, we fix the top-left square with a red chip. We apply casework to its two adjacent chips:

 <u>\textbf{Case (1): The top-center and center-left chips have different colors.}</u>
                    \begin{figure}[H]
                        \centering
                        \includegraphics[width=0.5\textwidth]{Solution/AMC_10_10A_2021_25_Solution_Figure_1.png}
                    \end{figure}
                    There are three subcases for Case (1):
                    \begin{figure}[H]
                        \centering
                        \includegraphics[width=0.5\textwidth]{Solution/AMC_10_10A_2021_25_Solution_Figure_2.png}
                    \end{figure}
                    As there are $3!=6$ permutations of the three colors, each subcase has $6$ ways. So, Case (1) has $3\cdot6=18$ ways in total.

 <u>\textbf{Case (2): The top-center and center-left chips have the same color.}</u>
                    \begin{figure}[H]
                        \centering
                        \includegraphics[width=0.5\textwidth]{Solution/AMC_10_10A_2021_25_Solution_Figure_3.png}
                    \end{figure}
                    There are three subcases for Case (2):
                    \begin{figure}[H]
                        \centering
                        \includegraphics[width=0.5\textwidth]{Solution/AMC_10_10A_2021_25_Solution_Figure_4.png}
                    \end{figure}
                    As there are $3!=6$ permutations of the three colors, each subcase has $6$ ways. So, Case (2) has $3\cdot6=18$ ways in total.

 <u>\textbf{Answer}</u>

 Together, the answer is $18+18=\boxed{\textbf{(E)} ~36}.$




\textbf{Solution 3 (Casework on the Red Chips' Configurations)}

We consider all possible configurations of the red chips for which rotations matter:
                    \begin{figure}[H]
                        \centering
                        \includegraphics[width=0.5\textwidth]{Solution/AMC_10_10A_2021_25_Solution_Figure_5.png}
                    \end{figure}
                    As there are $2!=2$ permutations of blue and green for each configuration, the answer is $2\cdot(2+4+4+4+4)=\boxed{\textbf{(E)} ~36}.$




\textbf{Solution 4 (Focusing On Top 2 Rows)}


                    \begin{figure}[H]
                        \centering
                        \includegraphics[width=0.5\textwidth]{Solution/AMC_10_10A_2021_25_Solution_Figure_6.png}
                    \end{figure}
                     

 We will first count the colorings where if the top two rows are filled out, the grid is uniquely determined. This will only happen if there are $3$ cells of one color, $2$ cells of another color, and $1$ cell of the remaining color in the top $3$ x $2$ grid made up of cells $1-6$. With this, we are left with two tokens of one color and one token of another color for the bottom row, so there must only be $1$ way to fill it up since the two tokens of the same color can't be next to each other. This makes the grid uniquely determined.

 The $3$ cells with the same color can either be cells $1,5,3$ or $2,4,6.$  If we choose $1,3,5,$ the two cells with the same color must be cells $4$ and $6$ and cell $2$ must be the cell of the other color. If we choose $2,4,6,$ the two cells with the same color must be cells $1$ and $3,$ and cell $5$ must be the cell of the other color. So there are $2$ ways to choose which cells share $3$ colors and $1$ way to choose which cells share $2$ and $1$ color. There are $3! = 6$ ways to assign colors to the cells. We have a total of $2 \cdot 1 \cdot 3! = 12$ colorings for this case. An example coloring is shown below with cells $1,3,5$ being chosen for the $3$ cells of the same color. 

 
                    \begin{figure}[H]
                        \centering
                        \includegraphics[width=0.5\textwidth]{Solution/AMC_10_10A_2021_25_Solution_Figure_7.png}
                    \end{figure}
                    

 Now, what about the colorings where even if the top two rows are filled out, the grid is still not uniquely determined?  In other words, what about the colorings where even if we fill in the top two rows, there is still more than one possible color combination for the last row?

 In this case, the $3$ x $2$ grid made up of the top two rows would have to use two tokens of each color, and the second row would be some permutation of the colors \textit{blue, red, green}. Then for the last row, we would be left with one token of each color. 

 There are $3! = 6$ ways to permute the second row. WLOG, say we fill the second row of the grid so that cell $4$ is red, cell $5$ is blue, and cell $6$ is green. Then cells $1,2,3$ can be \textit{green, red, blue} or \textit{blue, green, red}, respectively. Finally, cells $7,8,9$ can also either be \textit{green, red, blue} or \textit{blue, green, red}, respectively. This gives us $6 \cdot 4 = 24$ colorings. 

 
                    \begin{figure}[H]
                        \centering
                        \includegraphics[width=0.5\textwidth]{Solution/AMC_10_10A_2021_25_Solution_Figure_8.png}
                    \end{figure}
                    

 
Our final answer is $12 + 24 = \boxed{\textbf{(E)} ~36}$. 

 



\textbf{Solution 5 (Casework and Symmetry)}

$(1) \quad$$\begin{tabular}{ c c c }  R & G & ? \\   B & R & ? \\    ? & ? & ?     \end{tabular}$$\Longrightarrow$$\begin{tabular}{ c c c }  R & G & B \\   B & R & G \\    R & G & B     \end{tabular}$$\Longrightarrow$$\begin{tabular}{ c c c }  R & G & R \\   B & R & B \\    G & B & G     \end{tabular}$$\quad 3 \cdot 2 \cdot 2 = 12$

 There are $3$ choices for $R$, $2$ choices for $G$. $R$ on the down left corner can be switched with $B$ on the upper right corner.

 $(2) \quad$$\begin{tabular}{ c c c }  R & G & ? \\   B & R & ? \\    ? & ? & ?     \end{tabular}$$\Longrightarrow$$\begin{tabular}{ c c c }  R & G & B \\   B & R & G \\    G & B & R     \end{tabular}$$\quad 3 \cdot 2 = 6$

 There are $3$ choices for $R$, $2$ choices for $G$.

 $(3) \quad$$\begin{tabular}{ c c c }  G & R & ? \\   R & B & ? \\    ? & ? & ?     \end{tabular}$$\Longrightarrow$$\begin{tabular}{ c c c }  G & R & B \\   R & B & G \\    G & R & B     \end{tabular}$$\Longrightarrow$$\begin{tabular}{ c c c }  G & R & G \\   R & B & R \\    B & G & B     \end{tabular}$$\quad 3 \cdot 2 \cdot 2 = 12$

 Note that $(3)$ is a $180$° rotation of $(1)$.

 $(4) \quad$$\begin{tabular}{ c c c }  G & R & ? \\   R & B & ? \\    ? & ? & ?     \end{tabular}$$\Longrightarrow$$\begin{tabular}{ c c c }  G & R & B \\   R & B & G \\    B & G & R     \end{tabular}$$\quad 3 \cdot 2 = 6$

 Note that $(4)$ is a $90$° rotation of $(2)$.

 Therefore, the answer is $2 \cdot (12 + 6) = \boxed{\textbf{(E)} ~36}$.

 



\textbf{Solution 6 (Casework and Derangements)}

\textbf{Case (1)}: We have a permutation of R, B, and G as all of the rows. There are $3!$ ways to rearrange these three colors. After finishing the first row, we move onto the second. Notice how the second row must be a derangement of the first one. By the derangement formula, $\frac{3!}{e} \approx 2$, so there are two possible permutations of the second row. (Note: You could have also found the number of derangements of PIE). Finally, there are $2$ possible permutations for the last row. Thus, there are $3!\cdot2\cdot2=24$ possibilities.

 \textbf{Case (2)}: All of the rows have two chips that are the same color and one that is different. There are obviously $3$ possible configurations for the first row, $2$ for the second, and $2$ for the third. Thus, there are $3\cdot2\cdot2=12$ possibilities.

 Therefore, our answer is $24+12=\boxed{\textbf{(E)} ~36}.$




\textbf{Solution 7 (Rush, Use if you have less than 5 minutes)}

Ignore the center piece. Notice that when you place 3 of the chips, there are $6$ ways, making it inevitable for $2$ ways left for the other 2, so $6$ multiply $2$ is $12$. Now, notice there are $3$ ways to place a center piece, so our final statement is $12\cdot3=\boxed{\textbf{(E)} ~36}.$




\textbf{Solution 8 (Casework Using the Center and Surrounding)}

There are 3 indistinguishable red chips, 3 indistinguishable blue chips, and 3 indistinguishable green chips. Let's label the red chips R, the blue chips B, and the green chips G.

 WLOG, let's assume that the chip in the middle of the array is R.
 \[\begin{tabular}{ c c c }  ? & ? & ? \\   ? & R & ? \\    ? & ? & ?     \end{tabular}\]For the convenience of our proof, we will replace the question marks as follows.
 \[\begin{tabular}{ c c c }  1 & 2 & 3 \\   4 & R & 6 \\    7 & 8 & 9     \end{tabular}\]Doing this helps us easily identify where to place a certain chip. For example, if I want to place a green chip in the top leftmost square, I'd say, "Place a green chip in position 1."

 Now that we've assumed the center square is occupied by a red (R) chip, notice that the chips in positions $2, 4, 6, and 8$ can only be blue or green. With the information we now know, we can break the problem into cases based off of the number of a certain color of chip in the positions $2, 4, 6, 8$.

 \textbf{Case 1 (2 consecutive blue, 2 consecutive green):} An example case to show you what this means is shown below.
 \[\begin{tabular}{ c c c }  1 & G & 3 \\   B & R & G \\    7 & B & 9     \end{tabular}\]There are $4$ ways to figure out where the two blue and two green chips are placed. Notice that after I've placed the two blue, two green, and one red chip, the rest of the chips are in fixed places (try it!). Now we only need to worry about what color the center chip is.

 There are 3 ways to choose the color for the center chip, and afterwards, there are only two colors left for the two kinds of chips in positions $2, 4, 6,$ and $8$, and there are two of each, so there is no need to worry about the color of those chips.

 In conclusion, there are $4*3 = 12$ ways to arrange the chips in this case.

 \textbf{Case 2 (3 blue, 1 green OR 3 green, 1 blue):} An example case to show you what this means is shown below.
 \[\begin{tabular}{ c c c }  1 & G & 3 \\   B & R & G \\    7 & G & 9     \end{tabular}\]WLOG, let's first assume that there are $3$ green chips in such a position, and only $1$ blue chip. There are $4$ ways to figure out where they are placed. Notice that after I've placed the three green, one blue, and one red chip, the rest of the chips are yet again in fixed places (try it!). Now we only need to worry about what color the center chip is and which color of chip has $3$ in the positions of $2, 4, 6, 8$.

 There are $3$ ways to choose the color of the center chip. After that is chosen, there are only two colors left for the two kinds of chips needed in the positions $2, 4, 6,$ and $8$. However, since this time there are a different number of each color of chip, there are now $2$ ways to assign colors to the chips.

 In conclusion, there are $4*3*2 = 24$ ways to arrange the chips in this case.

 Notice how there should be a \textbf{Case 3} with cases that look like some variation of this:
 \[\begin{tabular}{ c c c }  1 & G & 3 \\   B & R & B \\    7 & G & 9     \end{tabular}\]There isn't, because if such an arrangement were to happen, we wouldn't be able to arrange the grid following the problem's conditions. There would either be a blue chip next to a blue chip, and/or a green chip next to a green chip.

 $12+24+0 = 36$, so the answer is $\boxed{\textbf{(E)} ~36}$.

 



\textbf{Solution 9}

Imagine a $3\times3$ grid as stated in the problem. Next, since this grid is extremely symmetric, lets assume the center square is $R$ (red). We realize that since $R$ is the center, the remaining $2$ $R's$ can either go in the top left square and the top right square or in the top left square and the bottom right square. Lets look at the first case. Note that the remaining $2$ $R's$ don't have in the top right and top left squares. We can rotate the board $4$ times to find the same case. So, lets multiply the final value in this case by $4$. Anyways, lets find the number of ways to color the rest of the board. Upon examination we can conclude that the top center square, the bottom left square, and the bottom right square all have to be the same color, so it can either be $B$ (blue) or $G$ (green), so we have $2$ choices. Thus, for this case, we have $4 \times 2 = 8$ ways to color the grid. Now, it is time for case $2$. Similarly, we need to multiply this case by $2$ because we can rotate the board twice to get the same case. So, upon some examination, we find that the top right square, the bottom center square, and left center square must be all the same color, so we've got $2$ choices. So in total we have $2\times2=4$ ways to color the grid this way. So in total, we have $12$ ways, but remember, we have to multiply by $3$ because we let the center square be $R$, so our answer is $12\times3=36 \Longrightarrow \boxed{\textbf{(E)}}$.

 -jb2015007

 


\end{enumerate}